\documentclass[11pt]{article}
\usepackage[margin=2.7cm]{geometry}
\usepackage{amsmath,amssymb,amsthm,mathtools}
\usepackage{enumitem}
\usepackage[colorlinks=true]{hyperref}
\usepackage[capitalize,noabbrev]{cleveref}

% ------------------------- theorem environments -------------------------
\newtheorem{theorem}{Theorem}[section]
\newtheorem{lemma}[theorem]{Lemma}
\newtheorem{corollary}[theorem]{Corollary}
\newtheorem{proposition}[theorem]{Proposition}
\theoremstyle{remark}
\newtheorem{remark}[theorem]{Remark}

% ------------------------- macros -------------------------
\newcommand{\R}{\mathbb{R}}
\newcommand{\bu}{\boldsymbol{u}}
\newcommand{\bv}{\boldsymbol{v}}
\newcommand{\bw}{\boldsymbol{w}}
\newcommand{\ba}{\boldsymbol{a}}
\newcommand{\bff}{\boldsymbol{f}}
\newcommand{\bn}{\boldsymbol{n}}
\newcommand{\ViscProj}{\Pi^{\mathrm{DS}}}
\newcommand{\norm}[1]{\lVert #1 \rVert}
\newcommand{\normK}[1]{\lVert #1 \rVert_{K}}
\newcommand{\tnorm}[1]{\lvert\!\lvert\!\lvert\, #1 \,\rvert\!\rvert\!\rvert}
\newcommand{\tns}{\tau_{1,\mathrm{NS}}}
\newcommand{\Eint}[2]{\mathrm{E}_{\mathrm{int},#1}(#2)}
\newcommand{\Dah}{\mathrm{Da}_{h,K}}
\newcommand{\Reh}{\mathrm{Re}_{h,K}}
\newcommand{\smul}{\circledast}
\newcommand{\pmul}{\circ}
\newcommand{\Pone}{\Pi_{1}}
\newcommand{\Ptwo}{\Pi_{2}}
\newcommand{\Ponez}{\Pi_{1,0}}
\newcommand{\Poneo}{\Pi_{1}^{\perp}}
\newcommand{\Ptwoo}{\Pi_{2}^{\perp}}
\newcommand{\Xhz}{\mathcal{X}_{h0}}

\title{Stability and a priori error analysis of the OSGS variant of a stabilized
finite element method for incompressible, inertial flows in inhomogeneous
porous media\\[1ex]
\large A self-contained companion note}
\author{}
\date{\today}

\begin{document}
\maketitle

\begin{abstract}
This note proves stability and a priori convergence for the orthogonal-subgrid-scale
(OSGS) variant of the stabilized finite element method introduced in the companion
manuscript~\cite{companion} for the linearized, incompressible, inertial porous-media
flow equations with variable porosity. The analysis mirrors the treatment of the Oseen
problem by orthogonal subscales in~\cite{Codina2008}, transported to the porosity-weighted
operator, the Darcy reaction term, and the artificial-compressibility perturbation of the
companion formulation. The main result is an inf--sup stability estimate in a mesh-dependent
norm that is \emph{stronger} than the one used for the ASGS variant in~\cite{companion}
(it controls the full reactive term $\norm{\sigma^{1/2}\bu_h}$ rather than the damped
$\lVert\widetilde{\sigma}_\alpha^{1/2}\bu_h\rVert$), together with the a priori estimate
\[
  \tnorm{U - U_h}
  \;\le\;
  C \sum_{K}
  \frac{\alpha_K}{h_K}\,\bigl(1 + \Dah\bigr)^{1/2}
  \Bigl( \tau_{2,K}^{1/2}\,\Eint{K}{\bu} + \tau_{1,K}^{1/2}\,\Eint{K}{p} \Bigr),
  \qquad
  \Dah \coloneqq \frac{\sigma h_K^2}{\alpha_K \nu},
\]
i.e., the ASGS error function of~\cite{companion} multiplied by the square root of one
plus the local mesh Damk\"ohler number. The factor is asymptotically harmless
($1 + O(h_K^2)$ under refinement) but permits a pre-asymptotic excess of order
$\Dah^{1/2}$ in strongly reaction-dominated regimes, and the mechanism that produces
it---the orthogonal projection annihilates the reactive subscale that, in the ASGS
variant, renormalizes $\sigma$ down to $\widetilde{\sigma}_\alpha$---is precisely
consistent with the numerical observations reported in the companion manuscript.
A $\sigma$-robust, optimal-order $L^2$ velocity estimate is obtained as a corollary.
\end{abstract}

\section{Introduction and scope}
\label{sec:intro}

The companion manuscript~\cite{companion} introduces two variational-multiscale
stabilized finite element methods for stationary, incompressible, inertial flows in
inhomogeneous porous media: an algebraic subgrid-scale (ASGS) variant and an
orthogonal subgrid-scale (OSGS) variant. The numerical analysis carried out
in~\cite{companion} (stability and a priori convergence of the linearized problem)
covers the ASGS variant only; the OSGS variant is exercised at length in the numerical
campaign but left without theory. The purpose of this note is to close that gap for the
linearized problem, following the template that~\cite{Codina2008} established for the
Oseen equations: an inf--sup argument built on a specially designed test function, a
consistency estimate based on the best-approximation property of weighted $L^2$
projections, and interpolation estimates in a mesh-dependent norm. The note is
self-contained; all hypotheses, parameters and norms are restated, in the notation
of~\cite{companion} wherever the two overlap.

\paragraph{Main result.}
Under the hypotheses (A1)--(A8) of \cref{sec:hypotheses}---which add to the standing
hypotheses of the ASGS analysis of~\cite{companion} only a Codina--Blasco-type stability
condition for the orthogonal projection (assumption (A7), the exact analogue of
hypothesis H4 in~\cite{Codina2008}) and, for interpolation orders $\ge 2$, extra smoothness
of the porosity field---the OSGS solution $U_h = [\bu_h; p_h]$ of the linearized problem
satisfies
\begin{equation}
  \label{eq:MainResultIntro}
  \tnorm{U - U_h}
  \;\le\;
  C \sum_{K}
  \frac{\alpha_K}{h_K}\,\bigl(1 + \Dah\bigr)^{1/2}
  \Bigl( \tau_{2,K}^{1/2}\,\Eint{K}{\bu} + \tau_{1,K}^{1/2}\,\Eint{K}{p} \Bigr),
\end{equation}
where $\Eint{K}{\psi} \coloneqq h_K^{k_\psi+1}\norm{\psi}_{H^{k_\psi+1}(K)}$ as
in~\cite{companion}, $\Dah \coloneqq \sigma h_K^2/(\alpha_K\nu)$ is the elementwise mesh
Damk\"ohler number, and $\tnorm{\cdot}$ is the OSGS triple norm defined in
\cref{sec:norm}, which dominates, term by term, the triple norm in which the ASGS
estimate of~\cite{companion} is stated. In other words: the OSGS bound is the ASGS bound
of~\cite{companion} with the same error function, in a norm at least as strong, times
$(1+\Dah)^{1/2}$. The constant $C$ is independent of $h$, $\nu$, $\sigma$,
$\varepsilon$ and of the porosity field except through the constants stated in
\cref{sec:hypotheses}.

\paragraph{The $(1+\Dah)^{1/2}$ factor and its mechanism.}
The factor is not an artifact of lazy estimation; it has a clean structural origin,
discussed in detail in \cref{sec:discussion}. In the ASGS variant, the
reactive--reactive coupling inside the residual-based stabilization term renormalizes
the reaction coefficient: the quadratic form contains
$\sigma(1-\sigma\tau_1)\norm{\bu_h}^2_K = \widetilde{\sigma}_\alpha\tau_1\cdot(\cdots)$,
and the damped coefficient $\widetilde{\sigma}_\alpha$ obeys
$\widetilde{\sigma}_\alpha \le c_1\alpha_K^2\tau_{2}/h^2$, so that reactive
contributions can always be traded against the divergence-stabilization scale without
any Damk\"ohler penalty. The orthogonal projection of the OSGS variant \emph{deletes}
the reactive part of the residual identically ($\sigma\bu_h$ is a finite element
function, hence has no fluctuation), so no renormalization takes place: the full
$\norm{\sigma^{1/2}\bu_h}^2$ survives on the left-hand side---a genuinely stronger
stability estimate---but the reactive interpolation terms must then be bounded head-on,
and the sharpest available conversion costs exactly $\Dah^{1/2}$. Both sides of this
trade are visible in the numerical campaign of~\cite{companion}: at strong reaction the
OSGS $L^2$ velocity error is indistinguishable from the ASGS one (the strong left-hand
side at work), while the OSGS $H^1$ velocity error carries an excess that shrinks, at
an accelerating relative rate, under refinement---the signature of a penalty
proportional to $\Dah^{1/2} \propto h$. \Cref{sec:discussion} makes this comparison
quantitative, with the caveats that one-sided upper bounds demand.

\paragraph{What is reused from the ASGS analysis.}
The proof reuses, unchanged: the Galerkin coercivity identity of~\cite{companion}
(the skew-symmetry of the convective and pressure--divergence pairings under
$\nabla\cdot(\alpha\ba)=0$ and Dirichlet conditions); the stabilization parameters
$\tau_1,\tau_2$ of~\cite{companion} and the elementary inequalities they satisfy
(collected in \cref{lem:parameters}, including one apparently new identity,
$\alpha_K^2\tau_1\tau_2 \ge h_K^2/(c_1+\Dah)$, which is the sole source of the
Damk\"ohler factor); the porosity-resolution condition and the upper bound on
$\varepsilon$; and the interpolation-error bookkeeping
$\Eint{K}{\cdot}$ together with the exact format of the ASGS convergence theorem, so
that the two results are directly comparable. Two hypotheses of the ASGS analysis are
\emph{not} needed here: the interior-face jump condition on $\tau_1$ and the enlarged
inverse-estimate constant for $\nabla\cdot(\alpha\nu\ViscProj\nabla\bu_h)$, because
the analyzed OSGS form contains no viscous or reactive stabilization terms to integrate
by parts (\cref{rem:dispensed}).

\paragraph{Honest scoping.}
As in~\cite{Codina2008}, the analyzed method is a first-order-consistent simplification
of the implemented one: the orthogonal projection is applied to the first-order part of
the residual, the adjoint test slot is truncated to its first-order part, and the
weighted inner product of the projection is the $\boldsymbol{\tau}$-weighted one rather
than the plain $L^2$ product used in practice. \Cref{rem:analyzed} spells out the exact
distance between the analyzed and the implemented formulations, item by item; the
annihilation lemma (\cref{lem:annihilation}) shows that for constant $\sigma$ and
$\varepsilon$ the reactive and compressibility contributions to both slots vanish
identically under the orthogonal projection, so that the only genuinely neglected
pieces are the viscous ones and the fluctuation of the force---precisely the Method~I
simplification of~\cite{Codina2008}. The projection stability condition (A7) is assumed,
not proved; \cite{Codina2008} indicates how the argument of~\cite{CodinaBlasco1997}
verifies its unweighted counterpart for simplicial elements. Finally, the analysis is
linear (given convection field $\ba$ with $\nabla\cdot(\alpha\ba)=0$, converged
projection), and it does not explain the \emph{superior} $L^2$ pressure convergence of
the OSGS variant observed for $k\ge 2$ in~\cite{companion}, which would require a
duality argument and remains open (\cref{rem:open}).

\section{The linearized problem and the OSGS method}
\label{sec:setting}

\subsection{Continuous problem and Galerkin form}
\label{sec:continuous}

Let $\Omega \subset \R^d$, $d = 2,3$, be a bounded polyhedral domain. Given a porosity
field $\alpha$, a solenoidal-in-the-weighted-sense convection field $\ba$ (see (A3),
(A4) below), constants $\nu > 0$, $\sigma \ge 0$, $\varepsilon \ge 0$, and a force
$\bff$, the linearized problem of~\cite{companion} reads: find $\bu$ and $p$ such that
\begin{subequations}
\label{eq:StrongProblem}
\begin{alignat}{2}
  \alpha\,\ba\cdot\nabla\bu
  \;-\; 2\,\nabla\cdot\bigl(\alpha\nu\,\ViscProj\nabla\bu\bigr)
  \;+\; \alpha\nabla p
  \;+\; \sigma\bu
  &= \bff
  &\qquad& \text{in } \Omega,
  \label{eq:StrongMomentum}\\
  \varepsilon p \;+\; \nabla\cdot(\alpha\bu) &= 0
  && \text{in } \Omega,
  \label{eq:StrongMass}\\
  \bu &= \boldsymbol{0} && \text{on } \partial\Omega,
  \label{eq:StrongBC}
\end{alignat}
\end{subequations}
where $\ViscProj$ is the pointwise orthogonal projection of second-order tensors onto
their deviatoric--symmetric part, as in~\cite{companion}. All boundary conditions are of
Dirichlet type, matching the setting of the ASGS analysis in~\cite{companion}. We write
$(\cdot,\cdot)$ and $\norm{\cdot}$ for the $L^2(\Omega)$ inner product and norm
(of scalars, vectors or tensors), $\normK{\cdot} \coloneqq \norm{\cdot}_{L^2(K)}$, and
$\norm{\cdot}_{m,K} \coloneqq \norm{\cdot}_{H^m(K)}$. The velocity space is
$V_0 \coloneqq H_0^1(\Omega)^d$ and the pressure space $Q_0$ is the zero-mean subspace
of $L^2(\Omega)$ (for $\varepsilon > 0$ the zero mean of the solution is automatic, and
whether or not the constraint is imposed on the space is immaterial to what follows;
see \cref{rem:meanshift}). Set $\mathcal{X}_0 \coloneqq V_0\times Q_0$ and, for
$U = [\bu; p]$, $V = [\bv; q]$,
\begin{equation}
  \label{eq:GalerkinForm}
  B(\ba; U, V)
  \coloneqq
  (\bv, \alpha\,\ba\cdot\nabla\bu)
  + 2(\nabla\bv, \alpha\nu\,\ViscProj\nabla\bu)
  + (\bv, \alpha\nabla p)
  + \sigma(\bv, \bu)
  + \varepsilon(q, p)
  + (q, \nabla\cdot(\alpha\bu)),
\end{equation}
which is the bilinear form of~\cite{companion}. Note that
$(\nabla\bv, \alpha\nu\ViscProj\nabla\bu) = (\ViscProj\nabla\bv, \alpha\nu\ViscProj\nabla\bu)$
because $\ViscProj$ is a pointwise orthogonal projection and $\alpha\nu$ is scalar.
The composite first-order quantity
\begin{equation}
  \label{eq:Xdef}
  X(V) \coloneqq \ba\cdot\nabla\bv + \nabla q
\end{equation}
plays throughout the role that $\ba\cdot\nabla\bv_h + \nabla q_h$ plays
in~\cite{Codina2008}; note that in the porous setting the momentum residual carries the
weight $\alpha$, i.e., the natural stabilized quantity is $\alpha X(V)$.

\begin{lemma}[Galerkin coercivity identity; reused from~\cite{companion}]
\label{lem:coercivity}
Assume $\nabla\cdot(\alpha\ba) = 0$ and let $U \in \mathcal{X}_0$ with
$\bu \in H^1_0(\Omega)^d$. Then
\begin{equation}
  \label{eq:CoercivityIdentity}
  B(\ba; U, U)
  = 2\nu\,\bigl\lVert \alpha^{1/2}\ViscProj\nabla\bu \bigr\rVert^2
  + \norm{\sigma^{1/2}\bu}^2
  + \varepsilon\norm{p}^2 .
\end{equation}
\end{lemma}

\begin{proof}
Integration by parts gives $(\bu, \alpha\nabla p) = -(p, \nabla\cdot(\alpha\bu))$, so
the pressure--divergence pair cancels, and
$(\bu, \alpha\ba\cdot\nabla\bu) = \tfrac12\bigl(\alpha\ba, \nabla|\bu|^2\bigr)
= -\tfrac12\bigl(\nabla\cdot(\alpha\ba), |\bu|^2\bigr) = 0$, both times using
$\bu|_{\partial\Omega} = \boldsymbol{0}$.
\end{proof}

\subsection{Meshes, finite element spaces and interpolation}
\label{sec:meshes}

Let $\mathcal{F} = \{\mathcal{T}_h\}_{h>0}$ be a family of conforming finite element
partitions of $\Omega$, $h_K \coloneqq \operatorname{diam}(K)$,
$h \coloneqq \max_K h_K$. We assume $\mathcal{F}$ nondegenerate (shape regular), so
that the inverse estimate
\begin{equation}
  \label{eq:InverseEstimate}
  \norm{\nabla v_h}_{K} \le \frac{C_{\mathrm{inv}}}{h_K}\norm{v_h}_{K}
\end{equation}
holds for all finite element functions and all $K$, and so that the element patches
$S_K$ (the union of elements sharing at least a vertex with $K$) are quasi-uniform:
$\underline{C}\, h_{K} \le h_{K'} \le \overline{C}\, h_K$ for all $K' \subset S_K$
\cite{AinsworthOden}. On $\mathcal{T}_h$ we build continuous finite element spaces
$V_h \subset H^1(\Omega)^d$ of degree $k_u \ge 1$ for the velocity and
$Q_h \subset H^1(\Omega)\cap C^0(\overline\Omega)$ of degree $k_p \ge 1$ for the
pressure (the pressure interpolation is continuous, as in~\cite{companion}), with
constrained subspaces $V_{h0} \coloneqq V_h \cap H^1_0(\Omega)^d$ and
$Q_{h0} \coloneqq Q_h \cap Q_0$, and $\Xhz \coloneqq V_{h0}\times Q_{h0}$. We assume
constants belong to $Q_h$ and $k_p \le k_u + 1$ (satisfied by equal order and by
Taylor--Hood-type pairs).

Two interpolation operators are used. For $v \in H^{r}(\Omega)$ with $r \ge 2$ (hence
$v \in C^0(\overline\Omega)$ for $d \le 3$), the nodal interpolant $\hat v_h$
satisfies
\begin{equation}
  \label{eq:NodalInterp}
  \norm{v - \hat v_h}_{m,K} \le C_I\, h_K^{n-m}\norm{v}_{n,K},
  \qquad 0 \le m \le n \le \min(r,\, k+1),
\end{equation}
with $k$ the polynomial degree of the target space, and it preserves homogeneous
Dirichlet data. For $v \in H^{r}(\Omega)$ with $r \ge 1$ the Scott--Zhang
quasi-interpolant~\cite{ScottZhang} satisfies the same estimate with $\norm{v}_{n,K}$
replaced by $\norm{v}_{H^n(S_K)}$ and also preserves homogeneous boundary values.

Finally, as in~\cite[Lemma~1]{Codina2008}, nondegeneracy furnishes continuous element-size
surrogates: there exist $\chi_1(\cdot,h), \chi_2(\cdot,h) \in C^0(\overline\Omega)$ with
$\chi_1(\boldsymbol{x}_K, h) \le h_K \le \chi_2(\boldsymbol{x}_K, h)$ at the element
barycenters $\boldsymbol{x}_K$ and
$\sup_h\sup_{\boldsymbol{x}}\chi_2/\chi_1 \le \chi_0 < \infty$. Following
\cite{Codina2008}, the element length entering the stabilization parameters below is
taken as $h_{K} \coloneqq \chi_1(\boldsymbol{x}_K, h)$; this alters each parameter by a
bounded factor ($\le \chi_0^2$) relative to using the raw diameter and is immaterial to
every estimate, so we do not distinguish the two notationally.

\subsection{Stabilization parameters}
\label{sec:parameters}

The stabilization parameters are those of~\cite{companion}: with
$\alpha_K \coloneqq \sup_K \alpha$ and
$|\ba|_{\infty,K} \coloneqq \norm{\ba}_{L^\infty(K)}$,
\begin{equation}
  \label{eq:TauDefs}
  \tns^{-1} \coloneqq \frac{c_1\nu}{h_K^2} + \frac{c_2\,|\ba|_{\infty,K}}{h_K},
  \qquad
  \tau_{1,K} \coloneqq \frac{1}{\alpha_K\,\tns^{-1} + \sigma},
  \qquad
  \tau_{2,K} \coloneqq \frac{h_K^2}{c_1\,\alpha_K\,\tns},
\end{equation}
with design constants $c_1, c_2 > 0$ subject to (A6) below. Expanding the definition,
\begin{equation}
  \label{eq:Tau2Expanded}
  \tau_{2,K}
  = \frac{\nu}{\alpha_K}\Bigl( 1 + \frac{c_2}{c_1}\,\Reh \Bigr),
  \qquad
  \Reh \coloneqq \frac{|\ba|_{\infty,K}\, h_K}{\nu},
\end{equation}
so that $\tau_{2,K} \ge \nu/\alpha_K$ always. We will also use the elementwise mesh
Damk\"ohler number $\Dah \coloneqq \sigma h_K^2/(\alpha_K\nu)$. Weighted broken inner
products and norms are defined, for $i = 1,2$, by
\begin{equation}
  \label{eq:TauInner}
  (f, g)_{\tau_i} \coloneqq \sum_K \tau_{i,K}\,(f, g)_{K},
  \qquad
  \norm{f}_{\tau_i} \coloneqq (f,f)_{\tau_i}^{1/2},
\end{equation}
for functions that need only be square integrable elementwise; note
$\norm{f}_{\tau_i} = \lVert \tau_i^{1/2} f\rVert_h$ in the broken-norm notation
of~\cite{companion}.

\subsection{Projections and the analyzed OSGS method}
\label{sec:method}

Let $\Pone$ denote the $\tau_1$-weighted $L^2$ projection onto the
\emph{unconstrained} velocity space $V_h$ (acting componentwise), $\Ptwo$ the
$\tau_2$-weighted projection onto the unconstrained pressure space $Q_h$, and
$\Poneo \coloneqq I - \Pone$, $\Ptwoo \coloneqq I - \Ptwo$ the complementary
(fluctuation) operators. We will also need $\Ponez$, the $\tau_1$-weighted projection
onto the \emph{constrained} space $V_{h0}$, which enters only the stability proof
through the special test function, exactly as in~\cite{Codina2008}. Projecting onto the
unconstrained spaces in the method matches both the implementation described
in~\cite{companion} (footnote to the OSGS system there) and the discussion
in~\cite[Remark~1]{Codina2008}: projecting onto the constrained spaces would simplify
the analysis (assumption (A7) would not be needed) but produces spurious numerical
boundary layers.

The \emph{analyzed} OSGS method reads: find $U_h \in \Xhz$ such that
\begin{equation}
  \label{eq:Method}
  B_{\mathrm{osgs}}(\ba; U_h, V_h) = L(V_h) \coloneqq (\bff, \bv_h)
  \qquad \forall\, V_h \in \Xhz,
\end{equation}
with
\begin{equation}
  \label{eq:Bosgs}
  B_{\mathrm{osgs}}(\ba; U, V)
  \coloneqq
  B(\ba; U, V) + S_1(U, V) + S_2(U, V),
\end{equation}
\begin{equation}
  \label{eq:Sdefs}
  S_1(U, V) \coloneqq \bigl( \Poneo[\alpha X(U)],\, \alpha X(V) \bigr)_{\tau_1},
  \qquad
  S_2(U, V) \coloneqq \bigl( \Ptwoo[\nabla\cdot(\alpha\bu)],\, \nabla\cdot(\alpha\bv) \bigr)_{\tau_2}.
\end{equation}
Since $(\Poneo z, \Pone w)_{\tau_1} = 0$ for all $z, w$, both stabilization terms are
symmetric and positive semidefinite:
\begin{equation}
  \label{eq:Ssym}
  S_1(U, V) = \bigl( \Poneo[\alpha X(U)],\, \Poneo[\alpha X(V)] \bigr)_{\tau_1},
  \qquad
  S_2(U, V) = \bigl( \Ptwoo[\nabla\cdot(\alpha\bu)],\, \Ptwoo[\nabla\cdot(\alpha\bv)] \bigr)_{\tau_2}.
\end{equation}
In particular, with \cref{lem:coercivity},
\begin{equation}
  \label{eq:Diagonal}
  B_{\mathrm{osgs}}(\ba; U_h, U_h)
  = 2\nu\,\bigl\lVert \alpha^{1/2}\ViscProj\nabla\bu_h \bigr\rVert^2
  + \norm{\sigma^{1/2}\bu_h}^2
  + \varepsilon\norm{p_h}^2
  + \norm{\Poneo \xi_h}_{\tau_1}^2
  + \norm{\Ptwoo \delta_h}_{\tau_2}^2,
\end{equation}
where here and below we abbreviate
\begin{equation}
  \label{eq:XiDelta}
  \xi_h \coloneqq \alpha X(U_h) = \alpha\bigl(\ba\cdot\nabla\bu_h + \nabla p_h\bigr),
  \qquad
  \delta_h \coloneqq \nabla\cdot(\alpha\bu_h).
\end{equation}

The next elementary observation locates the analyzed method precisely relative to the
implemented one, in which the orthogonal projection is applied to the \emph{full}
residual and the test slot is the full adjoint.

\begin{lemma}[Exact annihilations]
\label{lem:annihilation}
Let $\sigma$ and $\varepsilon$ be constant. Then, for all $U_h, V_h \in \Xhz$:
\begin{enumerate}[label=\textup{(\roman*)}, itemsep=0.2ex]
  \item $\Poneo(\sigma\bu_h) = \boldsymbol{0}$ and $\Ptwoo(\varepsilon p_h) = 0$
        (reactive and compressibility parts of the residual have no fluctuation);
  \item $\bigl( \Poneo z,\, \sigma\bv_h \bigr)_{\tau_1} = 0$ and
        $\bigl( \Ptwoo w,\, \varepsilon q_h \bigr)_{\tau_2} = 0$ for any $z, w$
        (reactive and compressibility parts of the adjoint test slot are invisible to
        the fluctuations).
\end{enumerate}
Consequently, the formulation obtained from the implemented OSGS method by projecting
the full residual and testing with the full adjoint differs from
\cref{eq:Method}--\cref{eq:Sdefs} only through (a) the fluctuations of the viscous
residual term $2\nabla\cdot(\alpha\nu\ViscProj\nabla\bu_h)$ and of the force $\bff$,
and (b) the viscous part of the adjoint test slot---exactly the Method~I simplification
of~\cite{Codina2008}.
\end{lemma}

\begin{proof}
For (i), $\sigma\bu_h \in V_h$ and $\varepsilon p_h \in Q_h$, and the fluctuation of
any element of the projection space vanishes. For (ii), $\sigma\bv_h \in V_h$ and
$\varepsilon q_h \in Q_h$, and fluctuations are $\tau_i$-orthogonal to the projection
spaces. The final statement follows by expanding the implemented stabilization
term by linearity.
\end{proof}

\begin{remark}[Analyzed versus implemented formulation]
\label{rem:analyzed}
For later reference we list the exact distance between \cref{eq:Method} and the OSGS
method as implemented in~\cite{companion}. (i)~\emph{Weighted versus plain projection}:
the analysis uses the $\boldsymbol{\tau}$-weighted projections of
\cref{eq:TauInner}; the implementation uses the plain $L^2$ projection. The two
coincide when $\tau_i$ is globally constant; otherwise the weighted product is the one
for which the best-approximation property of \cref{lem:bestapprox} holds, which is
essential to the consistency estimate (cf.~\cite[Remark~4]{Codina2008}).
(ii)~\emph{First-order truncation}: residual and test slots are truncated to their
first-order parts; by \cref{lem:annihilation} the only genuinely neglected pieces are
the viscous fluctuation, the force fluctuation, and the viscous test term. For $k_u = 1$
and elementwise-constant $\alpha\nu$ the viscous residual vanishes on element interiors;
with variable $\alpha$ it contributes $2\nu(\ViscProj\nabla\bu_h)\nabla\alpha$, which is
small under the resolution condition of (A3). (iii)~\emph{Converged versus lagged
projection}: the analysis takes the projection equation solved exactly (the coupled
system), whereas the implementation lags the projection along the nonlinear iteration.
(iv)~\emph{Unconstrained projection space}: analyzed and implemented formulations
agree (both project onto $V_h\times Q_h$ without Dirichlet constraints); note that the
display defining the OSGS projection in~\cite{companion} shows the projection onto
$\Xhz$, while the accompanying footnote states the implemented choice, which is the one
analyzed here and the one recommended by~\cite[Remark~1]{Codina2008}.
(v)~\emph{Elementwise-constant versus pointwise parameters}: as for the ASGS analysis
of~\cite{companion}, $\tau_{i}$ is elementwise constant in the analysis and pointwise
variable in the computations.
\end{remark}

\section{Hypotheses}
\label{sec:hypotheses}

Throughout, $C$ denotes a generic positive constant, independent of $h$, $\nu$,
$\sigma$, $\varepsilon$ and of $\alpha$ and $\ba$ except through the constants named
below; its value may change between appearances. We collect the standing assumptions:

\begin{itemize}[leftmargin=3.2em, itemsep=0.35ex]
  \item[(A1)] \emph{Mesh.} The family $\mathcal{F}$ is nondegenerate; the inverse
  estimate \cref{eq:InverseEstimate} holds with constant $C_{\mathrm{inv}}$; patches
  are quasi-uniform; the element length in \cref{eq:TauDefs} is the continuous
  surrogate of \cref{sec:meshes}.
  \item[(A2)] \emph{Coefficients.} $\nu > 0$, $\sigma \ge 0$ and $\varepsilon \ge 0$
  are constant, and the Darcy tensor is isotropic,
  $\boldsymbol{\sigma} = \sigma\boldsymbol{1}$, as in the ASGS analysis
  of~\cite{companion}.
  \item[(A3)] \emph{Porosity.} $\alpha \in W^{1,\infty}(\Omega)$ with
  $0 < \alpha_0 \le \alpha \le \alpha_\infty$ on $\overline\Omega$, and the mesh
  resolves it: there is $C_\nabla \ge 0$ with
  $h_K \norm{\nabla\alpha}_{L^\infty(K)} \le C_\nabla\, \alpha_K$ for all $K$ and all
  $h$ (the resolution condition of~\cite{companion}). For the consistency estimate with
  $m \coloneqq \max(k_u, k_p) \ge 2$ we additionally assume
  $\alpha \in W^{m,\infty}(\Omega)$ and set
  $C_{\alpha,m} \coloneqq 1 + \sum_{j=1}^{m}\norm{D^j\alpha}_{L^\infty(\Omega)}/\alpha_0$;
  for $k_u = k_p = 1$ no smoothness beyond $W^{1,\infty}$ is used.
  \item[(A4)] \emph{Convection field.} $\ba \in C^0(\overline\Omega)^d$, elementwise
  $W^{1,\infty}$, with $\nabla\cdot(\alpha\ba) = 0$ in the weak sense, and (only for
  the consistency estimate, cf.\ hypothesis H2 of~\cite{Codina2008})
  $\norm{D^j\ba}_{L^\infty(K)} \le C_D\, |\ba|_{\infty,K}$ for $1 \le j \le k_u$ and
  all $K$.
  \item[(A5)] \emph{Spaces.} Continuous velocity and pressure spaces of degrees
  $k_u, k_p \ge 1$ with $k_p \le k_u + 1$; constants belong to $Q_h$; all boundary
  conditions of Dirichlet type.
  \item[(A6)] \emph{Design constants and compressibility.} There is $\gamma \ge \gamma_0$
  such that $c_1 \ge \gamma^2 C_{\mathrm{inv}}^2$ and
  $c_2 \ge \gamma\, C_{\mathrm{inv}}$, where $\gamma_0$ depends only on $\beta_0$ of
  (A7), on $C_\nabla/C_{\mathrm{inv}}$, and on $C_2$ below (an admissible explicit
  choice is given in the proof of \cref{th:stability}); moreover
  $\varepsilon\,\tau_{2,K} \le C_2 \le 1$ for all $K$, matching the upper bound on
  $\varepsilon$ assumed in~\cite{companion} (there with $C_2 < 1$; $C_2 \le 1$ suffices
  here).
  \item[(A7)] \emph{Projection stability.} There is $\beta_0 > 0$, independent of $h$,
  such that for all $V_h \in \Xhz$, with $z_h \coloneqq \alpha X(V_h)$,
  \begin{equation}
    \label{eq:Hstab}
    \norm{z_h}_{\tau_1}
    \le \frac{1}{\beta_0}
    \Bigl( \norm{\Ponez z_h}_{\tau_1} + \norm{\Poneo z_h}_{\tau_1} \Bigr).
  \end{equation}
  This is the exact analogue of hypothesis H4 (condition (35)) of~\cite{Codina2008}:
  the component of $z_h$ in $V_h$ that is $\tau_1$-orthogonal to $V_{h0}$ is controlled
  by the other two. For the mass slot no analogous assumption is needed: the
  corresponding decomposition is exact with constant one (see Step~2 of the proof of
  \cref{th:stability}). As indicated in~\cite{Codina2008}, the analysis
  of~\cite{CodinaBlasco1997} can be used to verify the unweighted counterpart of
  \cref{eq:Hstab} for simplicial elements; we keep it as an assumption.
  \item[(A8)] \emph{Regularity.} The solution of \cref{eq:StrongProblem} satisfies
  $\bu \in H^{k_u+1}(\Omega)^d \cap H^1_0(\Omega)^d$ and
  $p \in H^{k_p+1}(\Omega) \cap Q_0$.
\end{itemize}

\section{Technical results}
\label{sec:technical}

\subsection{Parameter inequalities}

\begin{lemma}[Parameter inequalities]
\label{lem:parameters}
Under \textup{(A6)}, the parameters \cref{eq:TauDefs} satisfy, on every element $K$
(element subscripts omitted inside the formulas):
\begin{align}
  &\sigma\tau_1 \le 1,
  \qquad
  \alpha_K \tau_1 \le \tns,
  \qquad
  \varepsilon \le C_2\, \tau_2^{-1},
  \tag{P1}\label{eq:P1}\\
  &\tau_1\,\tau_2 \le \frac{h_K^2}{c_1\,\alpha_K^2},
  \tag{P2}\label{eq:P2}\\
  &h_K^2\,\tns^{-1} = c_1\,\alpha_K\,\tau_2,
  \qquad\text{hence}\qquad
  \bigl(\alpha_K\,\tns^{-1}\bigr)^{1/2} = \sqrt{c_1}\,\frac{\alpha_K}{h_K}\,\tau_2^{1/2},
  \tag{P3}\label{eq:P3}\\
  &\nu \le \alpha_K\,\tau_2,
  \qquad
  \tau_1^{1/2}\,|\ba|_{\infty,K} \le \frac{\sqrt{c_1}}{c_2}\,\tau_2^{1/2},
  \tag{P4}\label{eq:P4}\\
  &\alpha_K^2\,\tau_1\,\tau_2 \;\ge\; \frac{h_K^2}{c_1 + \Dah},
  \qquad\text{hence}\qquad
  \tau_2^{-1/2} \le \bigl(c_1 + \Dah\bigr)^{1/2}\,\frac{\alpha_K}{h_K}\,\tau_1^{1/2},
  \tag{P5}\label{eq:P5}\\
  &\sigma^{1/2} \;\le\; \Dah^{1/2}\,\frac{\alpha_K}{h_K}\,\tau_2^{1/2}.
  \tag{P6}\label{eq:P6}
\end{align}
Moreover, with $\gamma$ as in \textup{(A6)},
\begin{align}
  &C_{\mathrm{inv}}\,\frac{(\alpha_K\nu)^{1/2}}{h_K} \le \frac{1}{\gamma}\,\tau_1^{-1/2},
  \tag{K1}\label{eq:K1}\\
  &C_{\mathrm{inv}}\,\frac{\alpha_K\,|\ba|_{\infty,K}}{h_K} \le \frac{1}{\gamma}\,\tau_1^{-1},
  \tag{K2}\label{eq:K2}\\
  &C_{\mathrm{inv}}\,\frac{\alpha_K\,(\tau_1\tau_2)^{1/2}}{h_K} \le \frac{1}{\gamma}.
  \tag{K3}\label{eq:K3}
\end{align}
\end{lemma}

\begin{proof}
\cref{eq:P1}: $\tau_1^{-1} = \alpha_K\tns^{-1} + \sigma$ dominates each of its two
summands, giving the first two inequalities; the third is (A6). \cref{eq:P2}: by the
second inequality of \cref{eq:P1} and the definition of $\tau_2$,
$\tau_1\tau_2 \le (\tns/\alpha_K)\cdot h_K^2/(c_1\alpha_K\tns)$. \cref{eq:P3} is the
definition of $\tau_2$ rearranged. \cref{eq:P4}: the first inequality is
\cref{eq:Tau2Expanded}; for the second, $\tau_1 \le \tns/\alpha_K \le
h_K/(c_2\alpha_K|\ba|_{\infty,K})$ gives
$\tau_1|\ba|_{\infty,K}^2 \le |\ba|_{\infty,K}h_K/(c_2\alpha_K)$, while
\cref{eq:Tau2Expanded} gives
$\tau_2 \ge (c_2/c_1)\,|\ba|_{\infty,K}h_K/\alpha_K$; divide. \cref{eq:P5}: using
$\tns \le h_K^2/(c_1\nu)$,
\[
  \alpha_K^2\tau_1\tau_2
  = \frac{\alpha_K^2}{\alpha_K\tns^{-1}+\sigma}\cdot\frac{h_K^2}{c_1\alpha_K\tns}
  = \frac{h_K^2}{c_1\bigl(1 + \sigma\tns/\alpha_K\bigr)}
  \ \ge\ \frac{h_K^2}{c_1\bigl(1 + \sigma h_K^2/(c_1\alpha_K\nu)\bigr)}
  \ =\ \frac{h_K^2}{c_1 + \Dah}.
\]
\cref{eq:P6}: $\sigma^{1/2} = \Dah^{1/2}(\alpha_K\nu)^{1/2}/h_K \le
\Dah^{1/2}\alpha_K\tau_2^{1/2}/h_K$ by the first inequality of \cref{eq:P4}.
\cref{eq:K1}: $\tau_1^{-1} \ge \alpha_K\tns^{-1} \ge c_1\alpha_K\nu/h_K^2$, so
$C_{\mathrm{inv}}(\alpha_K\nu)^{1/2}/h_K \le
(C_{\mathrm{inv}}/\sqrt{c_1})\tau_1^{-1/2} \le \gamma^{-1}\tau_1^{-1/2}$ by (A6).
\cref{eq:K2}: $\tau_1^{-1} \ge \alpha_K\tns^{-1} \ge c_2\alpha_K|\ba|_{\infty,K}/h_K$
and $c_2 \ge \gamma C_{\mathrm{inv}}$. \cref{eq:K3}: combine \cref{eq:P2} with
$c_1 \ge \gamma^2C_{\mathrm{inv}}^2$.
\end{proof}

\subsection{Patch equivalence and the smoothing operation}

\begin{lemma}[Patch equivalence]
\label{lem:patch}
Under \textup{(A1)}, \textup{(A3)} and \textup{(A4)} there are constants
$0 < \underline{c} \le \overline{c}$, independent of $h$, such that for $i = 1,2$,
\begin{equation}
  \label{eq:PatchEquivalence}
  \underline{c}\;\tau_{i,K} \le \tau_{i,K'} \le \overline{c}\;\tau_{i,K},
  \qquad
  \underline{c}\;\alpha_{K} \le \alpha_{K'} \le \overline{c}\;\alpha_{K},
  \qquad \forall\, K' \subset S_K,\ \forall\, K.
\end{equation}
\end{lemma}

\begin{proof}
Patch quasi-uniformity gives $h_{K'} \eqsim h_K$. Since $\alpha$ is Lipschitz and
bounded below by $\alpha_0 > 0$, and $\operatorname{diam}(S_K) \le C h_K$, one has
$|\alpha_{K'} - \alpha_K| \le C h_K\norm{\nabla\alpha}_{L^\infty(S_K)} \le
C\,C_\nabla\,\max_{K''\subset S_K}\alpha_{K''}$, from which the uniform equivalence of
the $\alpha_{K}$ over the patch follows for $C\,C_\nabla$ bounded (and, in general, from
$\alpha_{K'}/\alpha_K \le \alpha_\infty/\alpha_0$; the sharper route through $C_\nabla$
keeps the constant independent of $\alpha_\infty/\alpha_0$). Likewise
$\ba \in C^0(\overline\Omega)$ implies $|\ba|_{\infty,K'} \le |\ba|_{\infty,K} +
\omega_{\ba}(Ch_K)$, where $\omega_{\ba}$ is the modulus of continuity, and the
resulting perturbations of $\tns^{-1}$, $\tau_1^{-1}$ and $\tau_2$ are controlled as
in the proof of \cref{lem:smoothing} below. The claim follows from the explicit
formulas \cref{eq:TauDefs}.
\end{proof}

Following~\cite{Codina2008}, for $\tau \in \{\tau_1, \tau_2\}$ and a finite element
function $v_h$ (scalar or vector) we define the discontinuous and the smoothed
products
\begin{equation}
  \label{eq:SmulDef}
  (\tau \pmul v_h)\big|_K \coloneqq \tau_{K}\, v_h\big|_K,
  \qquad
  (\tau \smul v_h)\big|_K \coloneqq \sum_{a \in N_K} N_a(\boldsymbol{x})\big|_K\,
  \bar\tau_{a}\, v_a,
\end{equation}
where $N_a$ are the nodal shape functions, $v_a$ the nodal values of $v_h$, and the
nodal parameters $\bar\tau_{i,a}$ are obtained by evaluating the formulas
\cref{eq:TauDefs} with $h_K$ replaced by $\chi_1(\boldsymbol{x}_a, h)$,
$\alpha_K$ by $\alpha(\boldsymbol{x}_a)$ and $|\ba|_{\infty,K}$ by
$|\ba(\boldsymbol{x}_a)|$. Thus $\tau \smul v_h$ is a continuous finite element
function, whereas $\tau \pmul v_h$ is not.

\begin{lemma}[Smoothing; adaptation of~{\cite[Lemma~2]{Codina2008}}]
\label{lem:smoothing}
Under \textup{(A1)}, \textup{(A3)}, \textup{(A4)} there is a function $\psi(h) \to 0$
as $h \to 0$ such that, for $\tau \in \{\tau_1, \tau_2\}$ and any finite element
function $v_h$,
\begin{equation}
  \label{eq:Smoothing}
  \norm{\tau\pmul v_h - \tau\smul v_h}_{K} \le \tau_K\,\psi(h)\,\normK{v_h},
  \qquad\text{hence}\qquad
  \norm{\tau\smul v_h}_{K} \le (1+\psi(h))\,\tau_K\,\normK{v_h}.
\end{equation}
If, in addition, $\ba$ is Lipschitz on $\overline\Omega$, then $\psi(h) = O(h)$.
\end{lemma}

\begin{proof}
As in~\cite[Lemma~2]{Codina2008}, the equivalence
$C_1 h_K^{d/2}\norm{w_h}_{L^\infty(K)} \le \normK{w_h} \le
C_2 h_K^{d/2}\norm{w_h}_{L^\infty(K)}$ for finite element functions on nondegenerate
meshes reduces \cref{eq:Smoothing} to the statement
$\max_{a\in N_K} |\bar\tau_a/\tau_K - 1| \le C\,\psi(h)$ with $\psi(h)\to 0$
uniformly in $K$. Write $\Delta(\cdot)$ for the difference between the nodal and the
elemental value of a quantity. For $\tau_1$,
$\tau_{1}|\Delta(\tau_1^{-1})| \le
\tau_1\bigl[\,|\Delta\alpha|\,\tns^{-1}
+ \alpha_\infty\, c_2\,|\Delta|\ba||/h_K
+ \alpha_K\, c_1\nu\,|\Delta(h^{-2})| + \alpha_K c_2 |\ba|_{\infty,K}|\Delta(h^{-1})|\,\bigr]$.
Using $\tau_1 \le \tns/\alpha_0$ and $\tau_1 \le h_K^2/(c_1\alpha_0\nu)$,
the first term is bounded by $\omega_\alpha(Ch)/\alpha_0$, the second by
$(c_2/(c_1\alpha_0\nu))\,h\,\omega_{\ba}(Ch)$, and the last two by
$C\,\omega_{\chi}(h)$, where $\omega_\alpha, \omega_{\ba}$ are the moduli of
continuity of $\alpha$ and $\ba$ and $\omega_\chi(h)\to 0$ quantifies the relative
oscillation of $\chi_1(\cdot,h)$ over an element, which vanishes as $h \to 0$ by the
construction of~\cite[Lemma~1]{Codina2008}. All right-hand sides tend to zero
uniformly; here the strict positivity $\alpha_0 > 0$ and $\nu > 0$ are used, and for
Lipschitz data all moduli are $O(h)$. Then
$|\bar\tau_{1,a}/\tau_{1,K} - 1| = \bar\tau_{1,a}\,|\Delta(\tau_1^{-1})|
\le (1+o(1))\,\tau_{1,K}|\Delta(\tau_1^{-1})|$, which proves the claim for $\tau_1$;
the same computation applied to the explicit form \cref{eq:Tau2Expanded} of $\tau_2$
(a continuous, positive expression in $\alpha$, $|\ba|$ and $h$, bounded away from
zero by $\nu/\alpha_\infty$) proves it for $\tau_2$.
\end{proof}

\subsection{Best approximation by the weighted projections}

\begin{lemma}[Best approximation and weighted residual approximation]
\label{lem:bestapprox}
Let $i \in \{1,2\}$ and let $\Pi_i$ be as in \cref{sec:method}. Then, for any
elementwise square-integrable $z$ and any finite element function $z_h$ in the
corresponding projection space,
\begin{equation}
  \label{eq:BestApprox}
  \norm{\Pi_i^{\perp} z}_{\tau_i} \le \norm{z - z_h}_{\tau_i}.
\end{equation}
Moreover, let $z \in H^{r}(\Omega)^s$ with $1 \le r \le k_u+1$ (for $i = 1$) or
$1 \le r \le k_p + 1$ (for $i = 2$), and let $\alpha$ satisfy \textup{(A3)} with
$m \ge r-1$ wherever $r - 1 \ge 2$. Then
\begin{equation}
  \label{eq:WeightedApprox}
  \norm{\Pi_i^{\perp}(\alpha z)}_{\tau_i}
  \;\le\;
  C\, C_{\alpha,r}
  \Bigl( \sum_K \tau_{i,K}\,\alpha_K^2\, h_K^{2r}\,\norm{z}_{H^{r}(K)}^2 \Bigr)^{1/2},
\end{equation}
with $C$ depending only on the interpolation constant, the patch-overlap constant and
the constants of \cref{lem:patch}, and $C_{\alpha,r}$ as in \textup{(A3)}
(with $C_{\alpha,1} \coloneqq 1 + \norm{\nabla\alpha}_{L^\infty(\Omega)}/\alpha_0$).
\end{lemma}

\begin{proof}
\cref{eq:BestApprox} is the minimizing property of orthogonal projections in their own
inner product, cf.~\cite[Remark~4]{Codina2008}. For \cref{eq:WeightedApprox}, take
$z_h$ to be the Scott--Zhang interpolant of $\alpha z$ in the projection space, so that
$\norm{\alpha z - z_h}_{K} \le C h_K^{r} \norm{\alpha z}_{H^{r}(S_K)}$
(here $r \le k+1$ for the relevant degree $k$, by (A5) and the stated restrictions on
$r$; in particular the case $i=1$, $z = \nabla p$, $r = k_p$ uses $k_p \le k_u+1$).
By the Leibniz rule and (A3),
$\norm{\alpha z}_{H^{r}(S_K)}
\le \bigl( \alpha_{\infty,S_K} + \sum_{j=1}^{r}\binom{r}{j}\norm{D^j\alpha}_{L^\infty}\bigr)
\norm{z}_{H^{r}(S_K)}
\le C\,C_{\alpha,r}\,\alpha_{\infty,S_K}\norm{z}_{H^{r}(S_K)}$,
using $\norm{D^j\alpha}_{L^\infty} \le (C_{\alpha,r}-1)\alpha_0 \le
(C_{\alpha,r}-1)\alpha_{\infty,S_K}$. Insert this into \cref{eq:BestApprox}, replace
$\alpha_{\infty,S_K}$ and $\tau_{i,S_K}$-values by their values on $K$ using
\cref{lem:patch}, and re-sum over patches using the bounded overlap implied by
nondegeneracy, absorbing $\norm{z}_{H^r(S_K)}$ into element contributions.
\end{proof}

\subsection{The OSGS triple norm}
\label{sec:norm}

The stability and convergence analysis is carried out in the mesh-dependent norm
\begin{equation}
  \label{eq:TripleNorm}
  \tnorm{V}
  \coloneqq
  \Bigl(
  2\nu\,\bigl\lVert \alpha^{1/2}\ViscProj\nabla\bv \bigr\rVert^2
  + \norm{\sigma^{1/2}\bv}^2
  + \varepsilon\norm{q}^2
  + \norm{\alpha X(V)}_{\tau_1}^2
  + \norm{\nabla\cdot(\alpha\bv)}_{\tau_2}^2
  \Bigr)^{1/2}.
\end{equation}

\begin{lemma}[Comparison with the ASGS norm]
\label{lem:normcomparison}
Let $\tnorm{\cdot}_{\mathrm{S}}$ denote the triple norm of the ASGS stability lemma
of~\cite{companion},
\[
  \tnorm{V}_{\mathrm{S}}^2
  = \nu\,\bigl\lVert \alpha^{1/2}\ViscProj\nabla\bv \bigr\rVert^2
  + \bigl\lVert \widetilde{\sigma}_\alpha^{1/2}\bv \bigr\rVert_h^2
  + \varepsilon\norm{q}^2
  + \norm{\alpha X(V)}_{\tau_1}^2
  + \norm{\nabla\cdot(\alpha\bv)}_{\tau_2}^2,
  \qquad
  \widetilde{\sigma}_\alpha
  = \frac{\tns^{-1}\sigma}{\tns^{-1} + \sigma/\alpha_K}.
\]
Then $\tnorm{V}_{\mathrm{S}} \le \tnorm{V}$ for all $V$.
\end{lemma}

\begin{proof}
Term by term: $\nu \le 2\nu$; on each element
$\widetilde{\sigma}_\alpha = \sigma\,\tns^{-1}/(\tns^{-1}+\sigma/\alpha_K) \le \sigma$;
the remaining three terms coincide.
\end{proof}

\section{Stability}
\label{sec:stability}

\begin{theorem}[Inf--sup stability]
\label{th:stability}
Under \textup{(A1)--(A7)} there exist $h_0 > 0$ and $\beta_{\mathrm{osgs}} > 0$,
depending only on $\beta_0$, $C_2$, $C_\nabla/C_{\mathrm{inv}}$ and the constants of
\cref{sec:technical}, such that for all $h \le h_0$,
\begin{equation}
  \label{eq:InfSup}
  \inf_{U_h \in \Xhz}\;
  \sup_{V_h \in \Xhz}\;
  \frac{B_{\mathrm{osgs}}(\ba; U_h, V_h)}{\tnorm{U_h}\,\tnorm{V_h}}
  \;\ge\; \beta_{\mathrm{osgs}} .
\end{equation}
An admissible choice of the constants in \textup{(A6)} is
$\gamma_0 = 10\,(1 + M_2)$ with $M_2 \coloneqq 1 + C_\nabla/C_{\mathrm{inv}}$,
together with $h_0$ small enough that
$\psi(h_0) \le \min\{1/8,\ \beta_0^2/16\}$, where $\psi$ is the function of
\cref{lem:smoothing}.
\end{theorem}

\begin{proof}
Fix $U_h = [\bu_h; p_h] \in \Xhz$ and recall the abbreviations \cref{eq:XiDelta}.
Write $\psi \coloneqq \psi(h)$ and set
\[
  D_\nu \coloneqq 2\nu\,\bigl\lVert \alpha^{1/2}\ViscProj\nabla\bu_h \bigr\rVert^2,
  \quad
  D_\sigma \coloneqq \norm{\sigma^{1/2}\bu_h}^2,
  \quad
  D_\varepsilon \coloneqq \varepsilon\norm{p_h}^2,
\]
\[
  O_1 \coloneqq \norm{\Poneo\xi_h}_{\tau_1}^2,
  \quad
  O_2 \coloneqq \norm{\Ptwoo\delta_h}_{\tau_2}^2,
  \quad
  P_1 \coloneqq \norm{\Ponez\xi_h}_{\tau_1}^2,
  \quad
  P_2 \coloneqq \norm{\Ptwo\delta_h}_{\tau_2}^2 .
\]

\medskip\noindent\emph{Step 1: the special test function.}
Define
\[
  \bw_h \coloneqq \tau_1 \smul \Ponez(\xi_h) \;\in\; V_{h0},
  \qquad
  r_h \coloneqq \tau_2 \smul \Ptwo(\delta_h) - m_h \;\in\; Q_{h0},
\]
where $m_h$ is the mean of $\tau_2\smul\Ptwo(\delta_h)$ over $\Omega$ whenever
$Q_{h0}$ consists of zero-mean functions, and $m_h \coloneqq 0$ otherwise. Note that
$\bw_h$ is an admissible velocity test function ($\Ponez\xi_h \in V_{h0}$ vanishes on
$\partial\Omega$, and the smoothed product preserves nodal zeros), and $r_h$ an
admissible pressure test function since constants belong to $Q_h$ by (A5). Set
$V^0 \coloneqq [\bw_h; r_h]$ and take as test function $W \coloneqq U_h + V^0$.
By \cref{eq:Diagonal},
\begin{equation}
  \label{eq:StepDiag}
  B_{\mathrm{osgs}}(\ba; U_h, U_h) = D_\nu + D_\sigma + D_\varepsilon + O_1 + O_2 ,
\end{equation}
and it remains to bound $B_{\mathrm{osgs}}(\ba; U_h, V^0)$ from below. Expanding
\cref{eq:GalerkinForm} and \cref{eq:Sdefs},
\begin{align}
  B_{\mathrm{osgs}}(\ba; U_h, V^0)
  ={}& 2\nu\,(\alpha\,\ViscProj\nabla\bu_h, \ViscProj\nabla\bw_h)
  \tag{T1}\label{eq:T1}\\
  &+ (\bw_h, \xi_h)
  \tag{T2}\label{eq:T2}\\
  &+ \sigma(\bu_h, \bw_h)
  \tag{T3}\label{eq:T3}\\
  &+ \varepsilon(p_h, r_h)
  \tag{T4}\label{eq:T4}\\
  &+ (r_h, \delta_h)
  \tag{T5}\label{eq:T5}\\
  &+ \bigl(\Poneo\xi_h,\ \alpha(\ba\cdot\nabla\bw_h + \nabla r_h)\bigr)_{\tau_1}
  \tag{T6}\label{eq:T6}\\
  &+ \bigl(\Ptwoo\delta_h,\ \nabla\cdot(\alpha\bw_h)\bigr)_{\tau_2},
  \tag{T7}\label{eq:T7}
\end{align}
where in \cref{eq:T2} we used
$(\bw_h, \alpha\ba\cdot\nabla\bu_h) + (\bw_h, \alpha\nabla p_h) = (\bw_h, \xi_h)$.
The constant shift $m_h$ is invisible throughout: it does not affect $\nabla r_h$ in
\cref{eq:T6}; in \cref{eq:T5} it contributes $-m_h\int_\Omega \delta_h = 0$ because
$\int_\Omega \nabla\cdot(\alpha\bu_h) = 0$ by the divergence theorem and
$\bu_h \in V_{h0}$; and in \cref{eq:T4} it contributes
$-\varepsilon\, m_h \int_\Omega p_h = 0$, since either $m_h = 0$ or $p_h$ has zero
mean. We now estimate \cref{eq:T1}--\cref{eq:T7} term by term. Throughout,
\cref{eq:Smoothing} is used in the form
$\normK{\bw_h} \le (1+\psi)\,\tau_{1,K}\normK{\Ponez\xi_h}$ and
$\normK{r_h + m_h} \le (1+\psi)\,\tau_{2,K}\normK{\Ptwo\delta_h}$ (recall also that
subtracting the mean does not increase the $L^2$ norm, so the same bound holds
globally for $r_h$).

\medskip\noindent\emph{Step 2: the two productive terms \cref{eq:T2} and \cref{eq:T5}.}
Splitting $\tau_1\smul(\cdot) = \tau_1\pmul(\cdot) + (\tau_1\smul - \tau_1\pmul)(\cdot)$
and using the $\tau_1$-orthogonality
$(\xi_h, \Ponez\xi_h)_{\tau_1} = \norm{\Ponez\xi_h}_{\tau_1}^2$,
\[
  (\bw_h, \xi_h)
  = P_1 + \bigl((\tau_1\smul - \tau_1\pmul)\Ponez\xi_h,\ \xi_h\bigr)
  \;\ge\; P_1 - \psi\,\norm{\Ponez\xi_h}_{\tau_1}\norm{\xi_h}_{\tau_1}
  \;\ge\; P_1 - \frac{2\psi}{\beta_0^2}\,(P_1 + O_1),
\]
where the last step used the nonexpansiveness of $\Ponez$ in $\norm{\cdot}_{\tau_1}$
together with \textup{(A7)}, in the form
$\norm{\xi_h}_{\tau_1}^2 \le (2/\beta_0^2)(P_1 + O_1)$. For \cref{eq:T5}, since
$\delta_h$ has zero mean the shift is invisible as noted, and by the exact
$\tau_2$-orthogonal decomposition of $\delta_h$ with respect to $Q_h$,
$\norm{\delta_h}_{\tau_2}^2 = P_2 + O_2$; hence
\[
  (r_h, \delta_h)
  = P_2 + \bigl((\tau_2\smul - \tau_2\pmul)\Ptwo\delta_h,\ \delta_h\bigr)
  \;\ge\; P_2 - \psi\,(P_2 + O_2).
\]
This is the point where the mass slot enjoys the exact analogue of \textup{(A7)} with
constant one: because $\delta_h$ has zero mean, projecting onto $Q_h$ and correcting
the test function by a constant replaces the constrained projection
$\Pi_{2,0}$ of~\cite{Codina2008} at no cost.

\medskip\noindent\emph{Step 3: the perturbation terms.}
For \cref{eq:T1}, by Cauchy--Schwarz, $|\alpha| \le \alpha_K$ elementwise,
$|\ViscProj \boldsymbol{M}| \le |\boldsymbol{M}|$ pointwise, the inverse estimate
\cref{eq:InverseEstimate}, \cref{eq:Smoothing} and \cref{eq:K1},
\begin{align*}
  \text{\cref{eq:T1}}
  &\ \ge\ -\,2\sum_K \nu^{1/2}\,\bigl\lVert\alpha^{1/2}\ViscProj\nabla\bu_h\bigr\rVert_K
  \cdot \nu^{1/2}\alpha_K^{1/2}\,\frac{C_{\mathrm{inv}}}{h_K}\,(1+\psi)\,\tau_{1,K}
  \norm{\Ponez\xi_h}_K\\
  &\ \ge\ -\,2 M_1 \Bigl(\frac{D_\nu}{2}\Bigr)^{1/2} P_1^{1/2},
\end{align*}
with $M_1 \coloneqq (1+\psi)/\gamma$, whence, by Young's inequality,
\[
  \text{\cref{eq:T1}} \ \ge\ -\frac{D_\nu}{4} - 2M_1^2\, P_1 .
\]
For \cref{eq:T3}, using $\sigma\tau_1 \le 1$ from \cref{eq:P1},
\begin{align*}
  \text{\cref{eq:T3}}
  &\ \ge\ -(1+\psi)\sum_K \sigma\,\tau_{1,K}\normK{\bu_h}\,\normK{\Ponez\xi_h}\\
  &\ \ge\ -(1+\psi)\, D_\sigma^{1/2} P_1^{1/2}
  \ \ge\ -\frac{D_\sigma}{2} - \frac{(1+\psi)^2}{2}\,P_1 .
\end{align*}
For \cref{eq:T4}, using $\varepsilon\tau_2 \le C_2$ from \cref{eq:P1},
\begin{align*}
  \text{\cref{eq:T4}}
  &\ \ge\ -(1+\psi)\sum_K \varepsilon\,\tau_{2,K}\normK{p_h}\normK{\Ptwo\delta_h}\\
  &\ \ge\ -(1+\psi)\,C_2^{1/2}\, D_\varepsilon^{1/2} P_2^{1/2}
  \ \ge\ -\frac{D_\varepsilon}{2} - \frac{(1+\psi)^2 C_2}{2}\,P_2 .
\end{align*}
For \cref{eq:T6} we bound the two contributions to
$\alpha X(V^0) = \alpha\ba\cdot\nabla\bw_h + \alpha\nabla r_h$ elementwise. By the
inverse estimate, \cref{eq:Smoothing} and \cref{eq:K2},
\[
  \tau_{1,K}^{1/2}\normK{\alpha\,\ba\cdot\nabla\bw_h}
  \le \tau_{1,K}^{1/2}\,\alpha_K |\ba|_{\infty,K}\frac{C_{\mathrm{inv}}}{h_K}
  (1+\psi)\,\tau_{1,K}\normK{\Ponez\xi_h}
  \le \frac{1+\psi}{\gamma}\,\tau_{1,K}^{1/2}\normK{\Ponez\xi_h},
\]
and by \cref{eq:K3},
\[
  \tau_{1,K}^{1/2}\normK{\alpha\nabla r_h}
  \le \tau_{1,K}^{1/2}\,\alpha_K\,\frac{C_{\mathrm{inv}}}{h_K}(1+\psi)\,
  \tau_{2,K}\normK{\Ptwo\delta_h}
  \le \frac{1+\psi}{\gamma}\,\tau_{2,K}^{1/2}\normK{\Ptwo\delta_h}.
\]
Hence, by Cauchy--Schwarz over the elements and $ab \le (a^2+b^2)/2$,
\[
  \text{\cref{eq:T6}}
  \ \ge\ -\frac{1+\psi}{\gamma}\, O_1^{1/2}\bigl( P_1^{1/2} + P_2^{1/2} \bigr)
  \ \ge\ -\frac{1+\psi}{\gamma}\Bigl( O_1 + \frac{P_1 + P_2}{2} \Bigr).
\]
For \cref{eq:T7}, expanding
$\nabla\cdot(\alpha\bw_h) = \alpha\,\nabla\cdot\bw_h + \nabla\alpha\cdot\bw_h$ and
using (A3) in the form
$\norm{\nabla\alpha}_{L^\infty(K)} \le C_\nabla\,\alpha_K/h_K$,
\[
  \tau_{2,K}^{1/2}\normK{\nabla\cdot(\alpha\bw_h)}
  \le \Bigl(1 + \frac{C_\nabla}{C_{\mathrm{inv}}}\Bigr)
  \tau_{2,K}^{1/2}\,\alpha_K\,\frac{C_{\mathrm{inv}}}{h_K}\,
  (1+\psi)\,\tau_{1,K}\normK{\Ponez\xi_h}
  \le \frac{M_2(1+\psi)}{\gamma}\,\tau_{1,K}^{1/2}\normK{\Ponez\xi_h},
\]
by \cref{eq:K3}, with $M_2 = 1 + C_\nabla/C_{\mathrm{inv}}$; hence
\[
  \text{\cref{eq:T7}}
  \ \ge\ -\frac{M_2(1+\psi)}{\gamma}\, O_2^{1/2} P_1^{1/2}
  \ \ge\ -\frac{M_2(1+\psi)}{2\gamma}\,\bigl( O_2 + P_1 \bigr).
\]

\medskip\noindent\emph{Step 4: collection.}
Adding \cref{eq:StepDiag} and the bounds of Steps~2--3,
\begin{align*}
  B_{\mathrm{osgs}}(\ba; U_h, W)
  \ \ge\ & \tfrac34\, D_\nu + \tfrac12\, D_\sigma + \tfrac12\, D_\varepsilon\\
  &+ \Bigl( 1 - \frac{2\psi}{\beta_0^2} - \frac{1+\psi}{\gamma} \Bigr)\, O_1
  + \Bigl( 1 - \psi - \frac{M_2(1+\psi)}{2\gamma} \Bigr)\, O_2\\
  &+ \Bigl( 1 - \frac{2\psi}{\beta_0^2} - 2M_1^2 - \frac{(1+\psi)^2}{2}
  - \frac{1+\psi}{2\gamma} - \frac{M_2(1+\psi)}{2\gamma} \Bigr)\, P_1\\
  &+ \Bigl( 1 - \psi - \frac{(1+\psi)^2 C_2}{2} - \frac{1+\psi}{2\gamma} \Bigr)\, P_2 .
\end{align*}
With $\gamma \ge \gamma_0 = 10(1+M_2)$ and
$\psi \le \psi_0 = \min\{1/8, \beta_0^2/16\}$ (so that $2\psi/\beta_0^2 \le 1/8$ and
$1+\psi \le 9/8$), a direct evaluation shows that every coefficient is at least
$1/8$ (for $P_1$: $1 - 0.125 - 0.007 - 0.633 - 0.029 - 0.057 \ge 0.14$; for $P_2$,
using $C_2 \le 1$: $1 - 0.125 - 0.633 - 0.029 \ge 0.21$; the others are larger).
Hence, with $C_0 = 1/8$,
\begin{equation}
  \label{eq:Collected}
  B_{\mathrm{osgs}}(\ba; U_h, W)
  \ \ge\ C_0\,\bigl( D_\nu + D_\sigma + D_\varepsilon + O_1 + O_2 + P_1 + P_2 \bigr).
\end{equation}

\medskip\noindent\emph{Step 5: recovery of the triple norm and boundedness of the test
function.}
By \textup{(A7)} and the exact decomposition of the mass slot,
\[
  \tnorm{U_h}^2
  = D_\nu + D_\sigma + D_\varepsilon + \norm{\xi_h}_{\tau_1}^2 + \norm{\delta_h}_{\tau_2}^2
  \le D_\nu + D_\sigma + D_\varepsilon + \frac{2}{\beta_0^2}(P_1 + O_1) + (P_2 + O_2),
\]
so \cref{eq:Collected} yields
$B_{\mathrm{osgs}}(\ba; U_h, W) \ge C_0 \min\{1, \beta_0^2/2\}\,\tnorm{U_h}^2$.
It remains to bound $\tnorm{W} \le \tnorm{U_h} + \tnorm{V^0}$. Using
\cref{eq:Smoothing} elementwise together with \cref{eq:K1} for the viscous term,
$\sigma\tau_1 \le 1$ for the reactive term, $\varepsilon\tau_2 \le C_2$ for the
compressibility term, and the bounds of Step~3 for
$\norm{\alpha X(V^0)}_{\tau_1}$ and $\norm{\nabla\cdot(\alpha\bw_h)}_{\tau_2}$,
\begin{align*}
  \tnorm{V^0}^2
  &\le \Bigl[ \frac{2(1+\psi)^2}{\gamma^2}
  + (1+\psi)^2
  + \frac{2(1+\psi)^2}{\gamma^2}
  + \frac{M_2^2(1+\psi)^2}{\gamma^2} \Bigr] P_1
  + \Bigl[ (1+\psi)^2 C_2 + \frac{2(1+\psi)^2}{\gamma^2} \Bigr] P_2\\
  &\le C\,(P_1 + P_2),
\end{align*}
and $P_1 \le \norm{\xi_h}_{\tau_1}^2$, $P_2 \le \norm{\delta_h}_{\tau_2}^2$ give
$\tnorm{V^0} \le C\,\tnorm{U_h}$, hence $\tnorm{W} \le C\,\tnorm{U_h}$. Combining,
\[
  \sup_{V_h\in\Xhz} \frac{B_{\mathrm{osgs}}(\ba; U_h, V_h)}{\tnorm{V_h}}
  \ \ge\ \frac{B_{\mathrm{osgs}}(\ba; U_h, W)}{\tnorm{W}}
  \ \ge\ \beta_{\mathrm{osgs}}\, \tnorm{U_h}
\]
with $\beta_{\mathrm{osgs}} = C_0\min\{1,\beta_0^2/2\}/C > 0$, which is
\cref{eq:InfSup}.
\end{proof}

\begin{remark}[On the structure of the proof]
\label{rem:proofstructure}
The proof is that of~\cite[Theorem~1]{Codina2008} with four modifications. (i)~All
inverse-estimate steps carry the porosity weight, absorbed through
\cref{eq:K1}--\cref{eq:K3}; only the plain inverse estimate
\cref{eq:InverseEstimate} is used---in contrast with the ASGS stability proof
of~\cite{companion}, no estimate for $\nabla\cdot(\alpha\nu\ViscProj\nabla\bv_h)$ is
needed, because the analyzed OSGS form contains no viscous stabilization term.
(ii)~The reactive pairing \cref{eq:T3} is new; it is absorbed by the \emph{full}
Galerkin reaction term $D_\sigma$ at the price of half of $P_1$, using nothing more
than $\sigma\tau_1 \le 1$. This is where the OSGS analysis retains
$\norm{\sigma^{1/2}\bu_h}$ in the norm, where the ASGS analysis retains only the
damped $\widetilde{\sigma}_\alpha$. (iii)~The compressibility pairing \cref{eq:T4} is
new and is controlled by the same condition $\varepsilon\tau_2 \le C_2$ assumed
in~\cite{companion}. (iv)~The constrained pressure projection $\Pi_{\tau_2,0}$
of~\cite{Codina2008} is replaced by the unconstrained projection plus a mean
correction of the test function, which is invisible in every term; this makes the
mass-slot stability constant exactly one and is the reason no analogue of
\textup{(A7)} is required for the divergence.
\end{remark}

\section{Convergence}
\label{sec:convergence}

Throughout this section, $U = [\bu; p]$ denotes the solution of
\cref{eq:StrongProblem} under the regularity (A8), and $U_h \in \Xhz$ the solution of
\cref{eq:Method}. Following~\cite{companion}, the local interpolation-error quantities
are
\begin{equation}
  \label{eq:EintDef}
  \Eint{K}{\bu} \coloneqq h_K^{k_u+1}\norm{\bu}_{H^{k_u+1}(K)},
  \qquad
  \Eint{K}{p} \coloneqq h_K^{k_p+1}\norm{p}_{H^{k_p+1}(K)},
\end{equation}
and the OSGS error function is defined by
\begin{equation}
  \label{eq:ErrorFunction}
  \mathcal{E}(h)
  \coloneqq
  \sum_K \frac{\alpha_K}{h_K}\,\bigl(1 + \Dah\bigr)^{1/2}
  \Bigl( \tau_{2,K}^{1/2}\,\Eint{K}{\bu} + \tau_{1,K}^{1/2}\,\Eint{K}{p} \Bigr),
\end{equation}
i.e., the error function of the ASGS convergence theorem of~\cite{companion} with the
local factor $(1+\Dah)^{1/2}$ inserted.

\subsection{Consistency}

The analyzed method is not consistent in the classical variational sense
(cf.~\cite[Eq.~(44)]{Codina2008}); its consistency error is the value of the
stabilization terms at the exact solution, and it is of the order of the error
function:

\begin{lemma}[Consistency error]
\label{lem:consistency}
Under \textup{(A1)--(A8)}, for all $V_h \in \Xhz$,
\begin{equation}
  \label{eq:ConsistencyId}
  B_{\mathrm{osgs}}(\ba;\, U - U_h,\, V_h)
  = S_1(U, V_h) + S_2(U, V_h),
\end{equation}
and
\begin{equation}
  \label{eq:ConsistencyBound}
  \bigl| B_{\mathrm{osgs}}(\ba;\, U - U_h,\, V_h) \bigr|
  \;\le\; C\,\mathcal{E}(h)\,\tnorm{V_h}.
\end{equation}
\end{lemma}

\begin{proof}
Since $U$ solves \cref{eq:StrongProblem} with the stated regularity, testing the
strong equations against $V_h \in \Xhz$ and integrating the viscous term by parts
gives $B(\ba; U, V_h) = L(V_h)$; adding $S_1(U,V_h) + S_2(U,V_h)$ to both sides and
subtracting \cref{eq:Method} yields \cref{eq:ConsistencyId}.

For the momentum part, split $\alpha X(U) = \alpha(\ba\cdot\nabla\bu) + \alpha\nabla p$
and use, respectively, Cauchy--Schwarz in $(\cdot,\cdot)_{\tau_1}$, linearity of
$\Poneo$ and \cref{lem:bestapprox} with $z = \ba\cdot\nabla\bu$, $r = k_u$ and with
$z = \nabla p$, $r = k_p$:
\[
  |S_1(U, V_h)|
  \le \Bigl( \norm{\Poneo[\alpha\,\ba\cdot\nabla\bu]}_{\tau_1}
  + \norm{\Poneo[\alpha\nabla p]}_{\tau_1} \Bigr)\,\norm{\alpha X(V_h)}_{\tau_1}.
\]
By (A4) and the Leibniz rule,
$\norm{\ba\cdot\nabla\bu}_{H^{k_u}(K)} \le C(1+C_D)\,|\ba|_{\infty,K}\,
\norm{\bu}_{H^{k_u+1}(K)}$, so \cref{eq:WeightedApprox} and the second inequality of
\cref{eq:P4} give
\[
  \norm{\Poneo[\alpha\,\ba\cdot\nabla\bu]}_{\tau_1}
  \le C \Bigl( \sum_K \tau_{1,K}\,\alpha_K^2\,|\ba|_{\infty,K}^2\, h_K^{2k_u}
  \norm{\bu}_{H^{k_u+1}(K)}^2 \Bigr)^{1/2}
  \le C \Bigl( \sum_K \Bigl(\frac{\alpha_K}{h_K}\Bigr)^{2} \tau_{2,K}\,
  \Eint{K}{\bu}^2 \Bigr)^{1/2},
\]
while $\norm{\nabla p}_{H^{k_p}(K)} \le \norm{p}_{H^{k_p+1}(K)}$ and
\cref{eq:WeightedApprox} give directly
\[
  \norm{\Poneo[\alpha\nabla p]}_{\tau_1}
  \le C \Bigl( \sum_K \tau_{1,K}\,\alpha_K^2\, h_K^{2k_p}
  \norm{p}_{H^{k_p+1}(K)}^2 \Bigr)^{1/2}
  = C \Bigl( \sum_K \Bigl(\frac{\alpha_K}{h_K}\Bigr)^{2} \tau_{1,K}\,
  \Eint{K}{p}^2 \Bigr)^{1/2}.
\]
Both right-hand sides are bounded by $C\,\mathcal{E}(h)$ (with room to spare: no
Damk\"ohler factor arises here), and
$\norm{\alpha X(V_h)}_{\tau_1} \le \tnorm{V_h}$.

For the mass part, the exact mass equation \cref{eq:StrongMass} gives
$\nabla\cdot(\alpha\bu) = -\varepsilon p$, hence, by \cref{lem:annihilation}-style
linearity,
\[
  S_2(U, V_h) = -\varepsilon\,\bigl( \Ptwoo p,\ \nabla\cdot(\alpha\bv_h)
  \bigr)_{\tau_2},
\]
which vanishes when $\varepsilon = 0$. For $\varepsilon > 0$, by Cauchy--Schwarz,
\cref{lem:bestapprox} (with $\alpha$ replaced by $1$, $z = p$, $r = k_p+1$),
$\varepsilon\tau_{2} \le C_2 \le 1$ and $\varepsilon \le C_2\tau_2^{-1}$ from
\cref{eq:P1}, and finally \cref{eq:P5},
\begin{align*}
  |S_2(U, V_h)|
  &\le \varepsilon \norm{\Ptwoo p}_{\tau_2}\,\norm{\nabla\cdot(\alpha\bv_h)}_{\tau_2}
  \le C \Bigl( \sum_K \varepsilon^2 \tau_{2,K}\, h_K^{2(k_p+1)}
  \norm{p}_{H^{k_p+1}(K)}^2 \Bigr)^{1/2} \tnorm{V_h}\\
  &\le C\,C_2^{3/2} \Bigl( \sum_K \tau_{2,K}^{-1}\, \Eint{K}{p}^2 \Bigr)^{1/2}
  \tnorm{V_h}\\
  &\le C \Bigl( \sum_K \bigl(c_1 + \Dah\bigr)
  \Bigl(\frac{\alpha_K}{h_K}\Bigr)^{2}\tau_{1,K}\,\Eint{K}{p}^2 \Bigr)^{1/2}
  \tnorm{V_h}
\end{align*}
(in the last step, $\varepsilon\tau_2 \le C_2 \le 1$ was applied once more, followed
by \cref{eq:P5}),
which is bounded by $C\,\mathcal{E}(h)\,\tnorm{V_h}$. Summing the two parts and using
$\bigl(\sum_K b_K^2\bigr)^{1/2} \le \sum_K b_K$ for nonnegative $b_K$ proves
\cref{eq:ConsistencyBound}.
\end{proof}

\subsection{Interpolation estimates}

Let $\hat U_h = [\hat\bu_h; \hat p_h]$ where $\hat\bu_h \in V_{h0}$ is the nodal
interpolant of $\bu$ (well defined since $\bu \in H^{k_u+1} \subset
C^0(\overline\Omega)$ for $d\le 3$, and boundary-value preserving) and
$\hat p_h \in Q_{h0}$ is the nodal interpolant of $p$, shifted by its mean when
$Q_{h0}$ consists of zero-mean functions. Write
$\hat E \coloneqq U - \hat U_h = [\hat{\boldsymbol e}_u; \hat e_p]$.

\begin{remark}[The pressure mean shift]
\label{rem:meanshift}
Since $\int_\Omega p = 0$, the shift constant is bounded by the raw interpolation
error, $|\Omega|^{1/2}|m| = |\int_\Omega (\hat p_h - p)|\,|\Omega|^{-1/2}
\le \norm{p - \hat p_h}$, so all interpolation estimates for $\hat e_p$ hold with at
most doubled constants. Moreover the constant part of $\hat e_p$ is invisible in the
pairings below: against $\nabla\cdot(\alpha\bv_h)$ it integrates to zero by the
divergence theorem, against $\varepsilon q_h$ it vanishes because $q_h$ has zero mean
(or no shift was performed), and it has no gradient. We therefore do not distinguish
the shifted interpolant notationally.
\end{remark}

\begin{lemma}[Interpolation continuity]
\label{lem:interpolation}
Under \textup{(A1)--(A8)}, for all $V_h \in \Xhz$,
\begin{equation}
  \label{eq:InterpolationBound}
  \bigl| B_{\mathrm{osgs}}(\ba;\, \hat E,\, V_h) \bigr|
  \;\le\; C\,\mathcal{E}(h)\,\tnorm{V_h}.
\end{equation}
\end{lemma}

\begin{proof}
Expand $B_{\mathrm{osgs}}(\ba; \hat E, V_h)$ into the Galerkin terms and the
stabilization terms and estimate each. Below, all interpolation estimates are
\cref{eq:NodalInterp} with $n = k_u+1$ or $n = k_p+1$ and $m \in \{0,1\}$.

\smallskip\noindent
\emph{(I1) Viscous term.} Using \cref{eq:P4} ($\nu \le \alpha_K\tau_{2,K}$),
\begin{align*}
  2\nu(\alpha\ViscProj\nabla\hat{\boldsymbol e}_u, \ViscProj\nabla\bv_h)
  &\le \sqrt{2}\,\tnorm{V_h}\Bigl( \sum_K 2\nu\alpha_K\,
  \norm{\nabla\hat{\boldsymbol e}_u}_K^2 \Bigr)^{1/2}\\
  &\le C\,\tnorm{V_h}\Bigl( \sum_K
  \Bigl(\frac{\alpha_K}{h_K}\Bigr)^2\tau_{2,K}\,\Eint{K}{\bu}^2 \Bigr)^{1/2}.
\end{align*}

\smallskip\noindent
\emph{(I2) Convective and mass--divergence terms, jointly.} Using
$\nabla\cdot(\alpha\ba) = 0$, $\hat{\boldsymbol e}_u \in H^1_0(\Omega)^d$
(nodal interpolation preserves the boundary values of $\bu$) and the continuity of the
pressure space, global integration by parts gives
\[
  (\bv_h, \alpha\ba\cdot\nabla\hat{\boldsymbol e}_u)
  + (q_h, \nabla\cdot(\alpha\hat{\boldsymbol e}_u))
  = -\bigl( \hat{\boldsymbol e}_u,\ \alpha\bigl(\ba\cdot\nabla\bv_h + \nabla q_h\bigr)
  \bigr)
  = -\bigl( \hat{\boldsymbol e}_u,\ \alpha X(V_h) \bigr).
\]
Hence, splitting $\tau_{1,K}^{-1/2} \le (\alpha_K\tns^{-1})^{1/2} + \sigma^{1/2}$ and
using \cref{eq:P3} for the first part and \cref{eq:P6} for the second,
\[
  \bigl|\bigl(\hat{\boldsymbol e}_u, \alpha X(V_h)\bigr)\bigr|
  \le \tnorm{V_h} \Bigl( \sum_K \tau_{1,K}^{-1} \norm{\hat{\boldsymbol e}_u}_K^2
  \Bigr)^{1/2}
  \le C\,\tnorm{V_h}\Bigl( \sum_K \bigl(c_1 + \Dah\bigr)
  \Bigl(\frac{\alpha_K}{h_K}\Bigr)^2 \tau_{2,K}\,\Eint{K}{\bu}^2 \Bigr)^{1/2}.
\]
This is the first of the two places where the Damk\"ohler factor enters.

\smallskip\noindent
\emph{(I3) Reactive term.} Using the full reactive control in the norm
\cref{eq:TripleNorm} and \cref{eq:P6},
\[
  \sigma(\hat{\boldsymbol e}_u, \bv_h)
  \le \norm{\sigma^{1/2}\bv_h}\Bigl( \sum_K \sigma\,
  \norm{\hat{\boldsymbol e}_u}_K^2 \Bigr)^{1/2}
  \le C\,\tnorm{V_h}\Bigl( \sum_K \Dah
  \Bigl(\frac{\alpha_K}{h_K}\Bigr)^2\tau_{2,K}\,\Eint{K}{\bu}^2 \Bigr)^{1/2}.
\]

\smallskip\noindent
\emph{(I4) Compressibility term.} Using $\varepsilon \le C_2\tau_2^{-1}$ and
\cref{eq:P5},
\[
  \varepsilon(\hat e_p, q_h)
  \le \varepsilon^{1/2}\norm{\varepsilon^{1/2}q_h}
  \Bigl( \sum_K \norm{\hat e_p}_K^2 \Bigr)^{1/2}
  \le C\,\tnorm{V_h}\Bigl( \sum_K \bigl(c_1+\Dah\bigr)
  \Bigl(\frac{\alpha_K}{h_K}\Bigr)^2\tau_{1,K}\,\Eint{K}{p}^2 \Bigr)^{1/2}.
\]

\smallskip\noindent
\emph{(I5) Pressure-gradient term.} Global integration by parts
($\bv_h \in V_{h0}$) and \cref{eq:P5} give
\[
  (\bv_h, \alpha\nabla\hat e_p)
  = -(\hat e_p, \nabla\cdot(\alpha\bv_h))
  \le \norm{\nabla\cdot(\alpha\bv_h)}_{\tau_2}
  \Bigl( \sum_K \tau_{2,K}^{-1}\norm{\hat e_p}_K^2 \Bigr)^{1/2}
  \le C\,\tnorm{V_h}\Bigl( \sum_K \bigl(c_1+\Dah\bigr)
  \Bigl(\frac{\alpha_K}{h_K}\Bigr)^2\tau_{1,K}\,\Eint{K}{p}^2 \Bigr)^{1/2},
\]
the second place where the Damk\"ohler factor enters (through
$\tau_2^{-1/2} \le (c_1+\Dah)^{1/2}(\alpha_K/h_K)\tau_1^{1/2}$).

\smallskip\noindent
\emph{(I6) Momentum stabilization.} Since fluctuation operators are nonexpansive,
\[
  |S_1(\hat E, V_h)|
  \le \bigl( \norm{\alpha\,\ba\cdot\nabla\hat{\boldsymbol e}_u}_{\tau_1}
  + \norm{\alpha\nabla\hat e_p}_{\tau_1} \bigr)\,\tnorm{V_h}.
\]
Elementwise, by \cref{eq:P4},
\[
  \tau_{1,K}^{1/2}\,\alpha_K\,|\ba|_{\infty,K}\norm{\nabla\hat{\boldsymbol e}_u}_K
  \le C\,\frac{\alpha_K}{h_K}\,\tau_{2,K}^{1/2}\,\Eint{K}{\bu},
  \qquad
  \tau_{1,K}^{1/2}\,\alpha_K\norm{\nabla\hat e_p}_K
  \le C\,\frac{\alpha_K}{h_K}\,\tau_{1,K}^{1/2}\,\Eint{K}{p},
\]
the second bound holding exactly; no Damk\"ohler factor arises.

\smallskip\noindent
\emph{(I7) Mass stabilization.} Likewise,
$|S_2(\hat E, V_h)| \le
\norm{\nabla\cdot(\alpha\hat{\boldsymbol e}_u)}_{\tau_2}\tnorm{V_h}$, and elementwise,
using (A3) in the form $\norm{\nabla\alpha}_{L^\infty(K)} \le C_\nabla\alpha_K/h_K$,
\[
  \tau_{2,K}^{1/2}\norm{\nabla\cdot(\alpha\hat{\boldsymbol e}_u)}_K
  \le \tau_{2,K}^{1/2}\bigl( \alpha_K \norm{\nabla\hat{\boldsymbol e}_u}_K
  + C_\nabla\tfrac{\alpha_K}{h_K}\norm{\hat{\boldsymbol e}_u}_K \bigr)
  \le C(1+C_\nabla)\,\frac{\alpha_K}{h_K}\,\tau_{2,K}^{1/2}\,\Eint{K}{\bu}.
\]

Collecting (I1)--(I7), bounding every square-summed right-hand side by the plain sum,
and absorbing $c_1$-dependent constants into $C$, we obtain
\cref{eq:InterpolationBound}.
\end{proof}

\begin{lemma}[Size of the interpolation error in the triple norm]
\label{lem:interpsize}
Under \textup{(A1)--(A8)},
$\tnorm{\hat E} \le C\,\mathcal{E}(h)$.
\end{lemma}

\begin{proof}
Estimate the five terms of \cref{eq:TripleNorm} at $V = \hat E$: the viscous term as
in (I1); $\norm{\sigma^{1/2}\hat{\boldsymbol e}_u}$ as in (I3);
$\varepsilon^{1/2}\norm{\hat e_p}$ as in (I4);
$\norm{\alpha X(\hat E)}_{\tau_1}$ as in (I6); and
$\norm{\nabla\cdot(\alpha\hat{\boldsymbol e}_u)}_{\tau_2}$ as in (I7).
\end{proof}

\subsection{Main theorem and corollaries}

\begin{theorem}[Convergence of the OSGS method]
\label{th:convergence}
Let $U$ solve the linearized problem \cref{eq:StrongProblem} and let
$U_h \in \Xhz$ solve \cref{eq:Method}. Under \textup{(A1)--(A8)} and for
$h \le h_0$ as in \cref{th:stability}, there is a constant $C$, independent of $h$,
$\nu$, $\sigma$, $\varepsilon$ and of $\alpha$, $\ba$ except through the constants of
\cref{sec:hypotheses}, such that
\begin{equation}
  \label{eq:MainResult}
  \tnorm{U - U_h}
  \;\le\;
  C \sum_K \frac{\alpha_K}{h_K}\,\bigl(1 + \Dah\bigr)^{1/2}
  \Bigl( \tau_{2,K}^{1/2}\,\Eint{K}{\bu} + \tau_{1,K}^{1/2}\,\Eint{K}{p} \Bigr).
\end{equation}
\end{theorem}

\begin{proof}
Write $U - U_h = \hat E + \eta_h$ with $\eta_h \coloneqq \hat U_h - U_h \in \Xhz$.
By \cref{th:stability} there is $V_h \in \Xhz$, $\tnorm{V_h} = 1$, with
\[
  \beta_{\mathrm{osgs}}\,\tnorm{\eta_h}
  \le B_{\mathrm{osgs}}(\ba; \eta_h, V_h)
  = -\,B_{\mathrm{osgs}}(\ba; \hat E, V_h)
  + B_{\mathrm{osgs}}(\ba; U - U_h, V_h)
  \le C\,\mathcal{E}(h),
\]
by \cref{lem:interpolation,lem:consistency}. The triangle inequality and
\cref{lem:interpsize} conclude.
\end{proof}

\begin{corollary}[Estimate in the ASGS norm]
\label{cor:asgsnorm}
Under the hypotheses of \cref{th:convergence}, the OSGS solution also satisfies the
estimate of the ASGS convergence theorem of~\cite{companion} up to the local factor
$(1+\Dah)^{1/2}$:
\[
  \tnorm{U - U_h}_{\mathrm{S}}
  \;\le\;
  C \sum_K \frac{\alpha_K}{h_K}\,\bigl(1 + \Dah\bigr)^{1/2}
  \Bigl( \tau_{2,K}^{1/2}\,\Eint{K}{\bu} + \tau_{1,K}^{1/2}\,\Eint{K}{p} \Bigr).
\]
\end{corollary}

\begin{proof}
Combine \cref{th:convergence} with \cref{lem:normcomparison}.
\end{proof}

\begin{corollary}[$\sigma$-robust $L^2$ velocity estimate]
\label{cor:Ltwo}
Let $\sigma > 0$. Under the hypotheses of \cref{th:convergence},
\begin{equation}
  \label{eq:LtwoVelocity}
  \norm{\bu - \bu_h}
  \;\le\;
  C \sum_K \bigl( 1 + \Dah^{-1} \bigr)^{1/2}
  \Bigl[
  \Bigl( 1 + \tfrac{c_2}{c_1}\Reh \Bigr)^{1/2} \Eint{K}{\bu}
  + \Bigl( \frac{\alpha_K}{\sigma\nu} \Bigr)^{1/2} \Eint{K}{p}
  \Bigr].
\end{equation}
In particular, wherever $\Dah \ge 1$ (the reaction-dominated pre-asymptotic regime),
the $L^2$ velocity error is of full interpolation order $h^{k_u+1}$ with a constant
uniform in $\sigma$, and the pressure contribution \emph{decreases} as
$(\sigma\nu)^{-1/2}$ as the reaction strengthens.
\end{corollary}

\begin{proof}
Since the norm \cref{eq:TripleNorm} controls $\norm{\sigma^{1/2}(\bu-\bu_h)}$,
$\norm{\bu - \bu_h} \le \sigma^{-1/2}\tnorm{U - U_h} \le
C\sigma^{-1/2}\mathcal{E}(h)$. For the velocity part of $\mathcal{E}(h)$, using
$\sigma^{-1/2} = \Dah^{-1/2}\,h_K/(\alpha_K\nu)^{1/2}$ and
\cref{eq:Tau2Expanded},
\[
  \sigma^{-1/2}(1+\Dah)^{1/2}\,\frac{\alpha_K}{h_K}\,\tau_{2,K}^{1/2}
  = (1+\Dah^{-1})^{1/2}\Bigl(\frac{\alpha_K}{\nu}\Bigr)^{1/2}\tau_{2,K}^{1/2}
  = (1+\Dah^{-1})^{1/2}\Bigl(1 + \tfrac{c_2}{c_1}\Reh\Bigr)^{1/2};
\]
for the pressure part, the same manipulation together with
$\tau_{1,K} \le 1/\sigma$ from \cref{eq:P1} gives
\[
  \sigma^{-1/2}(1+\Dah)^{1/2}\,\frac{\alpha_K}{h_K}\,\tau_{1,K}^{1/2}
  = (1+\Dah^{-1})^{1/2}\Bigl(\frac{\alpha_K}{\nu}\Bigr)^{1/2}\tau_{1,K}^{1/2}
  \le (1+\Dah^{-1})^{1/2}\Bigl(\frac{\alpha_K}{\sigma\nu}\Bigr)^{1/2}. \qedhere
\]
\end{proof}

\section{Discussion}
\label{sec:discussion}

\begin{remark}[The mechanism: $\sigma$ versus $\widetilde{\sigma}_\alpha$]
\label{rem:mechanism}
The single structural difference between the two analyses is the fate of the reactive
term. In the ASGS quadratic form the stabilization contains the reactive residual, and
the diagonal reactive contribution is
$\sigma\norm{\bu_h}_K^2 - \sigma^2\tau_{1}\norm{\bu_h}_K^2
= \sigma(1-\sigma\tau_1)\norm{\bu_h}_K^2$, with
$\sigma(1-\sigma\tau_1) = \sigma\,\alpha_K\tns^{-1}\tau_1 \eqsim
\widetilde{\sigma}_\alpha$: the subscales \emph{renormalize} the reaction down to
$\widetilde{\sigma}_\alpha \le \alpha_K\tns^{-1} = c_1\alpha_K^2\tau_2/h_K^2$, a
quantity always commensurate with the divergence-stabilization scale, so all reactive
pairings in the ASGS convergence proof can be traded against
$(\alpha_K/h_K)\tau_2^{1/2}$-weighted quantities with no Damk\"ohler penalty. In the
OSGS variant the orthogonal projection deletes the reactive residual identically
(\cref{lem:annihilation}), so no renormalization occurs: the diagonal keeps the full
$\norm{\sigma^{1/2}\bu_h}^2$ (a stronger left-hand side, cf.\
\cref{lem:normcomparison}), but the reactive and pressure interpolation pairings must
be bounded head-on, and the sharpest conversions available---\cref{eq:P5} and
\cref{eq:P6}---each cost a factor $(c_1+\Dah)^{1/2}$. The Damk\"ohler factor in
\cref{eq:MainResult} is thus the exact price of deleting the reactive subscale, and
the strengthened norm (in particular \cref{cor:Ltwo}) is the exact reward.
\end{remark}

\begin{remark}[Consistency with the numerical campaign of~\cite{companion}]
\label{rem:numerics}
The reaction-dominated experiments of~\cite{companion} (manufactured
solutions, $\mathrm{Da} = 10^6$) report two facts about the OSGS variant relative to the ASGS
one: (i)~the $L^2$ velocity errors of the two methods are indistinguishable at all
resolutions, and (ii)~the $H^1$-type velocity error of the OSGS method exceeds the
ASGS one by a factor of about $3.5$ on the finest $P_1$ meshes---the largest velocity
discrepancy of the whole campaign---with the gap closing under refinement at an
accelerating relative rate. Both observations sit exactly where the present analysis
puts them. Fact (i) is \cref{cor:Ltwo}: in the regime $\Dah \gg 1$ the OSGS $L^2$
velocity estimate is of full order with a $\sigma$-uniform constant, so no excess is
predicted (nor permitted at leading order) in $L^2$. Fact (ii) is consistent with the
factor $(1+\Dah)^{1/2}$ in \cref{eq:MainResult}: on the finest meshes of those
experiments the global mesh Damk\"ohler number is of order ten (its precise value
depends on the element-length convention entering $\tau$), so the bound permits an
excess of about $(1+\Dah)^{1/2} \approx 3$; on the interior plateau of the porosity
profile the elementwise values are larger and permit factors of $4$--$5$. Moreover,
$\Dah^{1/2} = (\sigma/(\alpha_K\nu))^{1/2}\, h_K$ decays \emph{linearly} in $h$ while
the bound itself decays like $h^{k_u}$, which reproduces the observed accelerating
relative decay of the gap. Two caveats are mandatory. First, upper bounds are
one-sided: the analysis shows that the observed excess is \emph{permitted} and
supplies a mechanism for it (\cref{rem:mechanism}); it does not prove that the excess
must occur, and the absence of the factor from the ASGS bound does not by itself prove
ASGS superiority. Second, the near-coincidence of the permitted and observed factors
at the finest resolution should be read as an order-of-magnitude consistency check,
not a validated prediction, since the constants in \cref{eq:MainResult} are not
tracked sharply.
\end{remark}

\begin{remark}[Hypotheses dispensed with]
\label{rem:dispensed}
Two technical hypotheses of the ASGS continuity analysis of~\cite{companion} are not
needed here: the interior-face jump condition on $\tau_1$ (required there when
$\sigma > 0$, following Eq.~(39) of~\cite{Codina2001}), and the enlarged
inverse-estimate constant for quantities of the form
$\nabla\cdot(\alpha\nu\ViscProj\nabla\bv_h)$. The reason is structural: the analyzed
OSGS form contains no viscous or reactive stabilization terms, so the proof never
integrates a stabilization term by parts elementwise---the only integrations by parts
are global ones, in (I2) and (I5) of \cref{lem:interpolation}. In exchange, the OSGS
analysis needs the projection stability condition \textup{(A7)}, which has no ASGS
counterpart. The two hypothesis sets are thus not nested; they reflect the different
routes (coercivity for ASGS, inf--sup for OSGS).
\end{remark}

\begin{remark}[A Method II variant]
\label{rem:methodII}
As in~\cite[\S 2.2]{Codina2008}, one may control the convective and pressure-gradient
fluctuations separately, replacing $S_1$ by
\[
  S_1^{\mathrm{II}}(U,V) \coloneqq
  \bigl( \Poneo[\alpha\,\ba\cdot\nabla\bu],\ \alpha\,\ba\cdot\nabla\bv \bigr)_{\tau_1}
  + \bigl( \Poneo[\alpha\nabla p],\ \alpha\nabla q \bigr)_{\tau_1}.
\]
The diagonal then controls the two fluctuations individually, and the same arguments
(with the elementary inequality $a^2 + b^2 \ge \tfrac13(a^2+b^2) + \tfrac13(a+b)^2$
used in~\cite[Theorem~3]{Codina2008}) give stability in the correspondingly finer
norm and the same convergence estimate \cref{eq:MainResult}. We have analyzed
Method~I because it is the variant implemented in~\cite{companion} (a single
projected quantity per equation).
\end{remark}

\begin{remark}[Robustness regimes and the screening length]
\label{rem:robustness}
Repeating the regime normalizations of the robustness section of~\cite{companion}
with \cref{eq:MainResult} in place of the ASGS estimate simply multiplies each regime
estimate by $(1+\Dah)^{1/2}$. In the viscous- and convection-dominated regimes
$\Dah \le 1$ on any reasonable mesh and the two variants obey identical bounds up to
constants; asymptotically, $\Dah \propto h_K^2$, so the OSGS bound exceeds the ASGS
one by $1 + O(h_K^2)$ only. The factor matters solely in the pre-asymptotic
reaction-dominated regime, and it has a clean physical reading: introducing the
porous screening length $\ell_\sigma \coloneqq (\alpha_K\nu/\sigma)^{1/2}$ (the decay
length of velocity perturbations in a Brinkman medium), one has
$\Dah = (h_K/\ell_\sigma)^2$, so the penalty persists exactly while the mesh
under-resolves $\ell_\sigma$ and disappears once $h_K \lesssim \ell_\sigma$. In the
$\mathrm{Da} = 10^6$ experiments of~\cite{companion}, $\ell_\sigma \sim 10^{-3}L$,
below the finest resolutions employed, which is why the excess is still visible there.
\end{remark}

\begin{remark}[What this analysis does not show]
\label{rem:open}
(i)~The numerical campaign of~\cite{companion} observes, for $k \ge 2$, an $L^2$
pressure convergence order for the OSGS variant one unit \emph{higher} than for the
ASGS variant. Nothing in the present analysis (nor in~\cite{Codina2008}, whose norms
contain no $L^2$ pressure term except in the viscous-dominated variant with
Reynolds-degrading constants) explains this; a duality argument exploiting the
orthogonality of the consistency error to the finite element space would be the
natural route, and remains open. (ii)~The analysis takes the projection equation
solved exactly; the lagged projection of the implemented staggered iteration, and the
nonlinear problem, are not covered. (iii)~The estimate \cref{eq:MainResult} is an
upper bound; no lower bound is claimed, so the comparison with the ASGS estimate in
\cref{cor:asgsnorm} orders the \emph{bounds}, not the methods. (iv)~For $\alpha \equiv 1$
and $\sigma = \varepsilon = 0$ the present theorem reduces to
\cite[Theorem~2]{Codina2008} with the parameter dictionary $c_3 = 1$,
$c_4 = c_2/c_1$, which serves as a consistency check of the constants.
\end{remark}

\subsection*{Acknowledgment of provenance}
This note was drafted as a working companion to~\cite{companion}; its numbering,
hypotheses and notation are meant to be lifted into that manuscript (or its
supplementary material) with minimal friction.

\begin{thebibliography}{10}

\bibitem{companion}
G.~Casas, J.~Gonz\'alez-Us\'ua, R.~Codina, I.~de-Pouplana,
\emph{A stabilized finite element method for incompressible, inertial flows in
inhomogeneous porous media}, companion manuscript (submitted).

\bibitem{Codina2008}
R.~Codina,
\emph{Analysis of a stabilized finite element approximation of the Oseen equations
using orthogonal subscales},
Applied Numerical Mathematics 58 (2008) 264--283.

\bibitem{Codina2001}
R.~Codina,
\emph{A stabilized finite element method for generalized stationary incompressible
flows},
Computer Methods in Applied Mechanics and Engineering 190 (2001) 2681--2706.

\bibitem{Codina2000}
R.~Codina,
\emph{Stabilization of incompressibility and convection through orthogonal sub-scales
in finite element methods},
Computer Methods in Applied Mechanics and Engineering 190 (2000) 1579--1599.

\bibitem{CodinaBlasco1997}
R.~Codina, J.~Blasco,
\emph{A finite element formulation for the Stokes problem allowing equal
velocity--pressure interpolation},
Computer Methods in Applied Mechanics and Engineering 143 (1997) 373--391.

\bibitem{Codina2002}
R.~Codina,
\emph{Stabilized finite element approximation of transient incompressible flows using
orthogonal subscales},
Computer Methods in Applied Mechanics and Engineering 191 (2002) 4295--4321.

\bibitem{BrennerScott}
S.~C. Brenner, L.~R. Scott,
\emph{The Mathematical Theory of Finite Element Methods},
Springer, New York, 2008.

\bibitem{ScottZhang}
L.~R. Scott, S.~Zhang,
\emph{Finite element interpolation of nonsmooth functions satisfying boundary
conditions},
Mathematics of Computation 54 (1990) 483--493.

\bibitem{HughesVMS}
T.~J.~R. Hughes, G.~R. Feij\'oo, L.~Mazzei, J.-B. Quincy,
\emph{The variational multiscale method---a paradigm for computational mechanics},
Computer Methods in Applied Mechanics and Engineering 166 (1998) 3--24.

\bibitem{AinsworthOden}
M.~Ainsworth, J.~T. Oden,
\emph{A Posteriori Error Estimation in Finite Element Analysis},
Wiley--Interscience, New York, 2000.

\end{thebibliography}

\end{document}
